\documentclass[11pt]{amsart}
\usepackage[T1]{fontenc}
\usepackage{lmodern,amsmath,amssymb,mathtools,microtype,booktabs}
\usepackage[margin=1.15in]{geometry}
\usepackage[hidelinks]{hyperref}
\newtheorem{theorem}{Theorem}[section]
\newtheorem{proposition}[theorem]{Proposition}
\newtheorem{lemma}[theorem]{Lemma}
\newtheorem{corollary}[theorem]{Corollary}
\theoremstyle{definition}\newtheorem{definition}[theorem]{Definition}
\theoremstyle{remark}\newtheorem{remark}[theorem]{Remark}
\newcommand{\R}{\mathbb R}
\newcommand{\Z}{\mathbb Z}
\newcommand{\E}{\mathbb E}
\newcommand{\Prob}{\mathbb P}
\newcommand{\one}{\mathbf 1}
\newcommand{\dd}{\,\mathrm d}
\DeclareMathOperator{\dist}{dist}
\DeclareMathOperator{\covol}{covol}
\DeclareMathOperator{\area}{area}
\DeclareMathOperator{\diam}{diam}
\numberwithin{equation}{section}
\title[Scalar collision laws]{Scalar collision laws and recognition of periodic dispersing billiards}
\author{Qian Qi}
\date{October 4, 2026}
\subjclass[2020]{37D50, 52A10, 60J10, 62G05, 68Q32}
\keywords{Dispersing billiards, reciprocal collision commands, active boundary estimation,
bit complexity, stationary launch noise, periodic recognition}
\hypersetup{pdftitle={Scalar collision laws and recognition of periodic dispersing billiards},pdfauthor={Qian Qi}}
\begin{document}
\raggedbottom
\begin{abstract}
We study geometric reconstruction of periodic dispersing tables from
reciprocal collision commands returning one bit per attempted preparation.
On bounded $C^{6,\beta}$ classes, write $s=6+\beta$. With shrinking
localization and controlled implementation errors, adaptive boundary
search attains $C^2$ accuracy $\nu$ with attempted-bit power
$\nu^{-1/(s-2)}$, matching deterministic and expected-length lower
powers up to logarithms. The commands have finite dyadic descriptions;
physical launch precision remains a separate resource. A converse at
position-error scale $\nu^{s/(s-2)}$ holds for two possibly different
table-, command- and draw-dependent bounded implementations, not for
one common jitter law. In contrast, we prove recovery with a fixed
calibrated convex launch support and an unknown stationary density
having a specified boundary-mass lower bound. For density bounded below
on that support, a finite reciprocal experiment has attempted-bit power
$\nu^{-(3s+1)/(s-2)}$ and no fixed-spread resolution floor. The inverse
recovers blurred occupation, identifies its Minkowski support components
and subtracts the known footprint, without learning the density. All
attempts, including solid starts and misses, are charged. Uniform finite
primitive-period decisions use a known positive nonperiod-patch margin;
real period vectors remain estimates. The results concern active
experiments with explicit apparatus priors, not passive counts or spectra.
\end{abstract}
\maketitle
\section{The observation and the main results}
\label{sec:setting}
The number of outcomes of a sensor and the information in its controls
are different resources. Here every preparation returns one bit, but
the controller chooses its spatial input. Our question is how many such
preparations are needed for geometric reconstruction, once their input
precision is specified. The reciprocal identity below reduces collision
observations to an occupation difference. The main point is that this
reduction can be evaluated at isolated, adaptively chosen locations;
a fine grid throughout the observation aperture is unnecessary.
We distinguish a localized experiment with worst-case implementation
errors from a stationary-noise experiment with a fixed calibrated launch
support. The latter, introduced below, has no shrinking random launch
spread and does not require the launch density to be known.

\subsection{Physical class and loss}
Fix $0<\beta\le1$ and put $s=6+\beta$. A table is the nonempty union
\begin{equation}\label{eq:presentation}
 \mathcal O=\bigcup_{i=1}^r(C_i+L\Z^2),\qquad 1\le r\le r_0.
\end{equation}
This presentation has exactly $r$ component orbits. It need not be
primitive. The constants in the following bounds are fixed and known:
\begin{equation}\label{eq:priors}
 \|L\|,\|L^{-1}\|\le L_0,\quad |\det L|\ge V_0>0,\quad
 \diam C_i\le D_0,\quad \dist(C,D)\ge d_0>0
\end{equation}
for distinct physical components. All components are strictly convex,
with curvature in $[\kappa_-,\kappa_+]\subset(0,\infty)$ and centered
support functions bounded by $M$ in $C^{6,\beta}(\R/2\pi\Z)$.
Reflections, repeated shapes and nonprimitive presentations are allowed.
Write $\mathcal B$ for this class, $\Pi=\{w:\mathcal O+w=\mathcal O\}$
for its full period group, $z_C$ for a component's Steiner point, and
$p_C$ for its centered support function. All norms on finite-dimensional
spaces are fixed once and for all.

A known positive nonperiod-patch margin $\eta$ will be required only
for uniform finite discrete decisions. Its precise definition is
\eqref{eq:patch-margin}. We denote the resulting subclass by
$\mathcal B_\eta$. The geometric loss is the maximum $C^2$ error of
matched component support functions in a protected laboratory patch,
together with the error of a matched primitive period basis and free
area. Representatives are chosen in a common protected patch. The
basis is matched to a true basis, not required to be a continuously
chosen reduced basis. Exact rational relations and the primitive orbit
count are stated separately from this real-valued loss.

\subsection{Commands and their accounting}
For a nominal displacement $a$, let $B_a(q)$ equal one if a free start
at $q$ has a first collision on the unreflected segment $[q,q+a]$.
It equals zero for a solid start and for a free miss. Both are attempted
preparations and enter the denominator. For pointwise statements solids and swept sets are closed: a boundary
solid start is zero, and a start on a swept boundary but outside every
solid is one. This fixes the tangency convention; changing it on null
sets has no effect on the means. Fix an even smooth probability kernel
$\kappa$ supported in the unit disk, and set
\[
 f_{x,\sigma}(q)=\sigma^{-2}\kappa((q-x)/\sigma),\qquad
 u_\sigma=\kappa_\sigma*\one_{\mathcal O}.
\]
Choose $\sigma_0,t>0$ with $t+2\sigma_0<d_0$, and prescribe the compass
law
\begin{equation}\label{eq:compass}
 \rho=\tfrac14(\delta_{te_1}+\delta_{-te_1}
                         +\delta_{te_2}+\delta_{-te_2}).
\end{equation}
A forward command draws $q\sim f_{x,\sigma}$ and $a\sim\rho$
independently and attempts $(q,a)$. A reverse command draws that same
nominal joint law and attempts $(q+a,-a)$. Separate fresh preparations
are used; pairing the same physical particle is unnecessary. Only the
two pooled means are used, not means resolved by direction. The joint
reverse law, rather than only its marginal direction law, is prescribed.

In this localized bounded-error experiment, each actual preparation may be coupled to its nominal one with position,
duration and angle errors bounded by $\ell,\tau,\alpha$, respectively.
It remains a unit-speed straight flight until its first impact. The
bounds hold conditionally on past commands, and preparations at a given
command are fresh and conditionally independent. Section~\ref{sec:queries}
proves the bias bound $C(\ell+\tau+t\alpha)/\sigma$. Calibration is a
resource of the controller, not an estimate inferred from the bits.
The launch aperture is fixed by the prior and does not grow as accuracy
increases. It includes the translated reverse starts and the kernel
supports. No free point, obstacle label, impact location, impact angle,
collision time, primitive cell or area normalization is supplied.

\begin{theorem}[Adaptive scalar reconstruction]\label{thm:adaptive}
Fix the numerical bounds of $\mathcal B_\eta$ and the kernel above.
There are $C,c,\nu_0>0$ such that, for $0<\nu<\nu_0$ and
$0<\delta<1/4$, a finite adaptive experiment with the reciprocal
compass commands reconstructs a smooth strictly convex periodic
presentation with geometric loss at most $C\nu$, with probability
at least $1-\delta$. It recovers the true primitive orbit count and
exact rational relations of a true period-generating list in a true
independent-pair basis. The Euclidean coordinates of these vectors
are estimates.

After a prior-dependent coarse acquisition, the finest localization
scale and sufficient coupled calibration are
\begin{equation}\label{eq:precision}
 \sigma=c\nu^{s/(s-2)},\qquad \ell+\tau+t\alpha\le c\sigma.
\end{equation}
The number of nominal command centers, counted with repetitions if
necessary, and the number of attempted preparations satisfy
\begin{align}
 J_\nu&\le C\nu^{-1/(s-2)}\log(C/\nu),\label{eq:centers}\\
 N_\nu&\le C\nu^{-1/(s-2)}\log(C/\nu)
                               \log(C/(\nu\delta)).\label{eq:attempts}
\end{align}
Every solid start and free miss is charged in $N_\nu$. All commands lie
in one fixed prior-controlled aperture. The reconstruction uses finite
local computations, radial bisections and finite polynomial data, not
a search over an infinite-dimensional class of physical candidates.
Its constants include the exit-time and patch-margin bounds.
\end{theorem}

Here and below a finite smooth partition of unity, with its derivatives,
is part of the fixed numerical representation. Certified evaluations
to the reserved tolerances suffice. The theorem counts preparations
and centers, not apparatus travel or bit operations for elementary
functions. Known $s$ is used; adaptivity refers to the locations of
commands, not adaptation to unknown smoothness.

\begin{theorem}[A necessary number of binary observations]
\label{thm:lower}
For each fixed $0<\beta\le1$ there is a fixed choice of the bounds
above and a nondegenerate subclass $\mathcal P\subset\mathcal B_\eta$
with the following property. Any adaptive randomized procedure whose
only table-dependent outputs are at most $N$ bits, and which recovers
the component boundaries with $C^2$ error at most $\nu$ with probability
at least $3/4$ uniformly on $\mathcal P$, must satisfy
\begin{equation}\label{eq:lower}
                         N\ge c\nu^{-1/(s-2)}.
\end{equation}
This holds even when the lattice, its primitive orbit count, component
correspondences and laboratory pose are given, and implementation is
exact. Input choices may have arbitrary precision and depend on all
previous bits. Independent controller randomness is allowed.
\end{theorem}

The construction of $\mathcal P$ is given in Section~\ref{sec:lower}.
The assertion is not a lower bound for every subclass, in particular
not for a singleton or a finite parametric family. On any bounded
class containing this construction, \eqref{eq:attempts} and
\eqref{eq:lower} identify the power of the attempted-bit complexity.
A logarithmic gap remains. Neither result prices input precision or
converts an active experiment into a passive billiard invariant.

\subsection{Finite controls, random stopping and physical resolution}
The next results add two resource converses to the preceding constructive
inverse. The geometric regularity remains $s=6+\beta$, and all uniform
finite global decisions still require the bounded periodic class and
known positive patch margin. No new geometric response is supplied.

\begin{theorem}[Attempted bits and calibration resolution]
\label{thm:resources}
The adaptive reciprocal experiment has a realization with dyadic centers
and scales, using $O(\log(C/\nu)+\log(C/\delta))$ binary digits per
batch command and the same attempted-bit bound
\eqref{eq:attempts}. On a fixed nondegenerate physical subclass, any
procedure whose only table-dependent data are collision bits and whose
stopping decision uses only its past bits and independent randomness
has worst-case expected attempt count at least
$c\nu^{-1/(s-2)}$ at uniform confidence $3/4$.

Moreover, under table-, command- and draw-dependent bounded-position-error
implementations,
uniform centered $C^2$ accuracy $\nu$ with confidence greater than
$1/2$ for short commands requires
\[
                            r\le C\nu^{s/(s-2)}.
\]
Here $r$ is the allowed position error, durations and directions may
be exact, and command lengths have any fixed bound less than the
component separation. Above a constant multiple of this scale,
two distinct physical tables admit identical complete adaptive bit
experiments under two possibly different admissible implementations.
No identical common-noise device is asserted. More attempted
bits do not remove this obstruction. The sufficient combined tolerance
$\ell+\tau+t\alpha\le c\nu^{s/(s-2)}$ remains as before.
\end{theorem}

Theorem~\ref{thm:digital} proves the finite-description claim and gives
a separate conservative numerical cost. Theorem~\ref{thm:expected}
proves the expected-stopping claim by a prefix-free transcript argument,
with the centering gauge explicit. Theorem~\ref{thm:calibration-floor}
proves the physical-resolution converse through a pointwise common-response
coupling of nearby convex swept bodies. Neither converse prices
apparatus production, and neither claims sharp logarithmic factors.
The deterministic-cap lower bound is retained independently: it is
not used as an unproved substitute for the sequential result.

\subsection{A contrasting experiment with fixed stationary jitter}
The preceding resolution converse does not apply to every apparatus
whose start has a bounded displacement. The next theorem recovers the
same table from a fixed, nonzero random launch spread. A known footprint
$K$ has positive curvature, a bounded $C^{6,\beta}$ support function and
$t+\diam K<d_0$. At every nominal center the apparatus draws a fresh
translate of one unknown density $j$, supported on $K$, with
$j(z)\ge b_0\dist(z,\partial K)^\gamma$ almost everywhere in its
interior. Here $b_0>0$ and $\gamma\ge0$ are known. The density is
stationary across all commands; it need not be symmetric or mean zero.
The controller does not observe the draw or its realized start.

\begin{theorem}[Fixed launch spread and unknown density]
\label{thm:fixed-spread-main}
On $\mathcal B_\eta$, under these footprint and density priors, one
finite reciprocal-compass controller, independent of $j$, reconstructs
the geometry to $C^2$ accuracy $C\nu$ and recovers the primitive discrete
data with probability at least $1-\delta$. It uses at most
\[
 C\nu^{-((2\gamma+3)s+1)/(s-2)}
               \log(C/\nu)\log(C/(\nu\delta))
\]
attempted bits in a fixed aperture. All solid starts and misses are
charged. The launch footprint is independent of $\nu$; only the nominal
command precision and probability-estimation accuracy are refined.
For a density bounded below on a fixed disk, $\gamma=0$ and the power
is $(3s+1)/(s-2)$. No optimality of this sample exponent is asserted.
\end{theorem}

Theorem~\ref{thm:stationary-jitter} proves the full statement.
The mechanism is geometric rather than Fourier deconvolution. The
reciprocal means recover a blurred occupation whose positive components
are $C+(-K)$. Subtraction of the known support function of $-K$ recovers
$C$. Lemma~\ref{lem:cap-mass} controls the small probabilities near
this expanded boundary. Proposition~\ref{prop:jitter-modulus} also
gives a H\"older inverse modulus uniformly over different unknown
stationary densities. Theorem~\ref{thm:calibration-floor}, by contrast,
uses two possibly different table-, command- and draw-dependent bounded
implementation maps. These are different uncertainty classes for the
same binary collision sensor, not conflicting universal resolution laws.

\section{From reciprocal collisions to local occupation queries}
\label{sec:queries}
This section supplies the observation reduction used by every adaptive
query. The random walk is computational: its exit time is not an
observed collision record or an apparatus stopping signal.

\begin{lemma}[Reciprocal balance and calibration]\label{lem:balance}
With $\chi=\one_{\mathcal O}$, outside null boundary sets,
\begin{equation}\label{eq:balance}
 B_a(q)-B_{-a}(q+a)=\chi(q+a)-\chi(q).
\end{equation}
Consequently the difference $g_\sigma$ of the two pooled nominal means
satisfies
\begin{equation}\label{eq:forcing}
 g_\sigma=(T-I)u_\sigma,\qquad
 Tv(x)=\tfrac14\sum_{a\in\{\pm te_1,\pm te_2\}}v(x+a).
\end{equation}
For either command, coupled errors of size
$r=\ell+\tau+t\alpha\le\sigma$ change its mean by at most $Cr/\sigma$.
The constant is uniform over the protected aperture, components and
nominal directions, including tangencies.
\end{lemma}
\begin{proof}
If both endpoints are free, a segment and its reverse either both hit
or both miss. If exactly one endpoint is solid, the direction starting
there gives zero and the other direction hits. If both are solid both
bits are zero. This proves \eqref{eq:balance}. Averaging in $q,a$ and
commuting translation with convolution proves \eqref{eq:forcing}.

For a convex body put $S_a(C)=C+[-a,0]$. The hit bit is
\[
 B_a(q)=\one_{\cup_C S_a(C)}(q)-\chi(q).
\]
The displacement error is at most $\tau+t\alpha$. Support inequalities
show that membership can change only if the nominal start lies in
an $r$-tube of $\partial C$ or $\partial S_a(C)$, including the position
error. The part of a convex boundary tube in a disk of radius $\sigma$
has area at most $C(\sigma+r)r$. One proof integrates boundaries of
outer parallel bodies and inner erosions by the coarea formula; their
portions in the disk have length at most $2\pi\sigma$, by monotonicity
of the perimeter of convex sets. Degenerate levels follow by a limit.
Only a bounded number of components contribute: each has a point in
a disk of radius $\sigma+t+r$, and such points from distinct components
are $d_0$-separated. Multiplication by the density bound
$\sigma^{-2}\|\kappa\|_\infty$ proves the assertion. The tube containment
is pointwise for every allowed error, so the coupling need not have
mean zero. An exceptional coupling event adds its probability.
\end{proof}

Put $D=D_0+2\sigma_0$, $H=(D+t)^2/t^2$ and
$L_* =\max(1,\lceil2H\rceil)$. The positive set of $u_\sigma$ lies in
the $\sigma$-enlarged components. They have diameter at most $D$ and
mutual distance greater than $t$.

\begin{lemma}[Stopped inverse in one aperture]\label{lem:stopped}
For the compass walk $X_n=x+a_1+\cdots+a_n$ and
$\tau_u=\inf\{n\ge0:u_\sigma(X_n)=0\}$ one has
\begin{equation}\label{eq:survival}
 \sup_x\E_x\tau_u\le H,\qquad
 \sup_x\Prob_x(\tau_u>n)\le2^{-\lfloor n/L_*\rfloor}.
\end{equation}
Let a bounded target region $U$ contain the complete protected bodies
and their search brackets. Choose a fixed bounded $W$ containing
$U+\overline B_{D+t+1}$. On $W$ use the killed transition
$T^Wv=T(\one_Wv)$ and set iterates to zero outside $W$. For
$V_0=0$, $V_{n+1}=\max(0,T^WV_n-g_\sigma)$,
\begin{equation}\label{eq:bellman-error}
 0\le u_\sigma-V_n\le2^{-\lfloor n/L_*\rfloor}\quad\hbox{on }U.
\end{equation}
Uniform forcing error $\xi$ and numerical update error $\zeta$ add
at most $n(\xi+\zeta)$ if approximate iterates are clipped to $[0,1]$.
\end{lemma}
\begin{proof}
Until its first zero, a path stays in one enlarged component, since
one step cannot reach another. For $S=\tau_u\wedge n$,
$|X_S-x|\le D+t$. The bounded stopping identity for the martingale
$|X_n-x|^2-nt^2$ gives $t^2\E S\le(D+t)^2$. Monotone convergence
proves the mean bound. Markov's inequality at $L_*$ and the Markov
property at successive blocks prove the geometric tail.

On $W$, $T^Wu_\sigma\le Tu_\sigma$, so $0\le V_n\le u_\sigma$
and the iterates vanish on its zero set. Starting in $U$, the complete
containing component and every first-exit location lie in $W$. Along
this stopped path the killed transition has no boundary discrepancy.
For $e_n=u_\sigma-V_n$, the recursion gives $e_{n+1}\le Te_n$ before
exit, and hence
\[
 e_n(x)\le\E_x[u_\sigma(X_n)\one_{\{\tau_u>n\}}].
\]
This proves \eqref{eq:bellman-error}. Maximum, averaging and clipping
are nonexpansive in supremum norm; induction proves the error bound.
No periodic wrap or renormalization of a boundary row is used.
\end{proof}

The stopped argument also compares two physical occupations with
unrelated zero sets. Indeed, for $w=u-\widetilde u$ and
$h=(T-I)w$, bounded stopping at $\tau_u\wedge n$ gives
\[
 w(x)=\E_xw(X_{\tau_u\wedge n})
                 -\E_x\sum_{j<\tau_u\wedge n}h(X_j).
\]
At exit the terminal value is $-\widetilde u\le0$. On survival it
is at most one. Letting $n$ increase and then interchanging $u$ and
$\widetilde u$ yields $\|u-\widetilde u\|_{U,\infty}
\le H\|(T-I)(u-\widetilde u)\|_{W,\infty}$. Thus exact reciprocal
functionals at all sufficiently small scales determine the local
occupation and, by approximate-identity convergence and regular
closedness, the physical set. This retained exact conclusion requires
neither analyticity nor a supplied zero set. The martingale and killed
Green mechanisms are classical; see \cite{LawlerLimic}.

\begin{proposition}[A finite local query from pooled bits]
\label{prop:query}
For any adaptively chosen $x\in U$ and $\sigma\le\sigma_0$, an
occupation estimate with error less than $1/8$ can be computed using
only a prior-bounded number of reciprocal command centers in $W$.
At conditional confidence $1-\varepsilon$ it uses at most
$C\log(C/\varepsilon)$ attempted preparations. A sufficient calibration
is $\ell+\tau+t\alpha\le c\sigma$. Thresholding the estimate at
$1/2$ gives the correct membership label outside the $\sigma$-tube
of the component boundaries. Its label inside that tube is unrestricted.
\end{proposition}
\begin{proof}
Fix $n_*=6L_*$. To evaluate $V_{n_*}(x)$, its dependency locations
are exactly among
\begin{equation}\label{eq:diamond}
 x+t(i,j),\qquad |i|+|j|\le n_*.
\end{equation}
Intersect these with $W$ and put zero at locations outside it. Evaluate
the recursion backwards through shrinking diamonds, initializing
$V_0=0$. Induction on the layer shows equality with the full killed
iterate at the target. There are $O(n_*^2)$ centers and $O(n_*^3)$
updates, both independent of $\sigma$. No fine background spatial grid,
interpolation of the transition, or growing aperture is involved.

At each center estimate each pooled mean to error at most
$1/(1024n_*)$, and impose the same bound on its calibration bias.
The forcing error is then at most $1/(256n_*)$. Reserve update error
at most $1/(256n_*)$. Lemma~\ref{lem:stopped} gives total error at most
$1/64+1/256+1/256<1/8$. Hoeffding's inequality and a union bound over
the fixed finite number of means give the preparation bound. The
constant includes $n_*$ and may be large. The proof holds conditional
on past data because the new attempts are fresh. Inside a component
at boundary distance greater than $\sigma$ the average is one;
outside all enlarged components it is zero. Thresholding therefore
has the stated property. A membership label is a computed conclusion
from these bits, not an extra response furnished by the sensor.
\end{proof}

The full launch region may be taken to contain
$W+\overline B_{t+\sigma_0+1}$, which includes reverse positions and
kernel supports. For at most $Q$ adaptive queries, allocating conditional
failure probability $\delta/Q$ to each gives joint correctness with
probability $1-\delta$. It suffices to know a deterministic upper bound
on $Q$; repeated centers may be resampled. No independence of the
adaptively selected locations themselves is asserted.

\section{Adaptive recovery of complete component boundaries}
\label{sec:boundary}
We first acquire a coarse, fixed-resolution description. Its purpose
is to find interior centers and isolated radial brackets, not to supply
the final accuracy. This avoids both a supplied free reference point
and a bisection that might accidentally encounter a different obstacle.

\begin{lemma}[Coarse hulls and interior centers]\label{lem:coarse}
A prior-dependent finite number of queries from
Proposition~\ref{prop:query}, at one fixed scale $\sigma_c>0$, recovers
all complete components needed in the protected patch to Hausdorff
error $e_c\le3\sigma_c$. It supplies, for each such component $C$, a
polygon $H$ and a computable center $o$ for which
\begin{equation}\label{eq:inner-ball}
 \overline B_r(o)\subset C\cap H\subset\overline B_{R_1}(o)
\end{equation}
with fixed positive $r,R_1$. For every direction $v$, it supplies a
radial interval of fixed length at most $C e_c$ containing the boundary
point of $C$ on $o+\R_+v$. The interval is disjoint from the
$\sigma$-enlargements of all other components for every sufficiently
small fine scale $\sigma$.
\end{lemma}
\begin{proof}
The curvature upper bound supplies an interior rolling radius
$r_* =1/\kappa_+$. In support coordinates this follows from
$p+p''\ge r_*$: subtracting the support function of the radius-$r_*$
disk leaves a support function with nonnegative curvature measure.
In particular $C=C_{-2\sigma_c}+2\sigma_c\overline B$ when
$2\sigma_c<r_*$. Choose $\sigma_c<\min(r_*/8,d_0/40)$ and a square
grid of covering radius $\rho_c\le\sigma_c/8$ on a surrounding
square with collar larger than $2D_0+2$.

At its nodes retain the estimated occupations at least $1/2$, join
retained nodes at distance at most $4\sigma_c$, and take a convex hull
for each connected cluster. Every retained node is within $\sigma_c$
of a unique body; all nodes of $C_{-\sigma_c}$ are retained. For any
retained node assigned to $C$, choose a point in $C_{-2\sigma_c}$ at
distance at most $3\sigma_c$, and then a nearest grid node. That node
is retained and joined to the first. Convexity connects any two such
deep nodes by a segment; partitioning it into lengths at most $\rho_c$
and replacing partition points by nearest nodes yields a retained
path with successive lengths at most $3\rho_c$. Different bodies
cannot join because $d_0-2\sigma_c>4\sigma_c$. The hull thus lies in
$C+\sigma_c\overline B$, whereas every point of $C$ is within
$2\sigma_c+\rho_c$ of a retained deep node. This proves the error.
Discard clusters whose hulls are within $D_0+1$ of the outer square
boundary. Fragments cannot survive this cutoff, and the chosen collar
protects every required complete component.

Compute a largest inscribed disk of $H$ by its finite half-plane
inequalities, to any fixed reserved precision. Since
$d_H(H,C)\le e_c$, the support inequalities imply
$C_{-e_c}\subset H$ and $H_{-e_c}\subset C$. A disk of radius $r_*$
in $C$ therefore gives an inscribed disk of radius at least $r_*-e_c$
in $H$. The computed center, with the radii reduced by the error and
numerical reserve, satisfies \eqref{eq:inner-ball} for a fixed $r>0$.
Diameter bounds supply $R_1$.

Write $R_H(v),R_C(v)$ for the radial functions about $o$. Since both
bodies contain $\overline B_r(o)$,
\[
 H+e_c\overline B\subset o+(1+e_c/r)(H-o),
\]
and the analogous inclusion holds for $C$. The two Hausdorff
inclusions give $|R_H(v)-R_C(v)|\le C e_c$ uniformly. The polygonal
radial value is computable from its sides. Use a slightly enlarged
bracket about it, allowing numerical error. By decreasing the fixed
$\sigma_c$, its length and every point's distance to $C$ are less
than $d_0/4$. Separation excludes every other component from the
bracket and its sufficiently small kernel neighborhood.
\end{proof}

\begin{lemma}[Bisection with an indeterminate boundary layer]
\label{lem:bisection}
On each bracket in Lemma~\ref{lem:coarse}, the derived label of
Proposition~\ref{prop:query} may be arbitrary, and may depend on the
past, only when
\begin{equation}\label{eq:radial-ambiguity}
                       |a-R_C(v)|\le A\sigma
\end{equation}
for a fixed $A$. After $k$ midpoint tests, ordinary bisection returns
an estimate with error at most $A\sigma+2^{-k-1}w_0$, where $w_0$
is the initial bracket length. In particular $O(\log(C/\sigma))$
queries recover $R_C(v)$ to $O(\sigma)$.
\end{lemma}
\begin{proof}
Translate $o$ to zero. The inner ball implies
\[
 (1-\sigma/r)C\subset C_{-\sigma},\qquad
 C+\sigma\overline B\subset(1+\sigma/r)C.
\]
Together with $R_C\le R_1$, these show that outside
\eqref{eq:radial-ambiguity} the label is the true radial inside/outside
label. Other components were excluded from the bracket. A label one
at radius $a$ implies $a\le R_C+A\sigma$; a label zero implies
$a\ge R_C-A\sigma$. If the current endpoints are $l,u$, midpoint
updates consequently preserve the relaxed invariant
$R_C\in[l-A\sigma,u+A\sigma]$. This does not require $R_C$ to remain
in $[l,u]$, or the noisy labels to be monotone. The width halves at
every step. Its midpoint has the asserted error.
\end{proof}

\begin{lemma}[Radial regularity and stable interpolation]
\label{lem:radial}
Under \eqref{eq:inner-ball} and the physical support priors, the
$2\pi$-periodic radial function $R_C(\varphi)$ has a uniformly bounded
$C^{6,\beta}$ norm. Let $h=2\pi/m$ for an integer $m\ge9$. From
radial values at the $m$ equally spaced angles, each with error at most
$e$, one can construct a smooth function $\widehat R$ such that
\begin{equation}\label{eq:radial-interpolation}
 \|\widehat R-R_C\|_{C^j}
           \le C(h^{s-j}+e h^{-j}),\qquad 0\le j\le3.
\end{equation}
For $e\le C_0h^s$ and sufficiently small $h$, the curve
$o+\widehat R(\varphi)(\cos\varphi,\sin\varphi)$ is the boundary of
a smooth strictly convex body $\widehat C$, and
\begin{equation}\label{eq:geometric-interpolation}
 d_H(\widehat C,C)\le Ch^s,\qquad
 \|p_{\widehat C}^{\rm lab}-p_C^{\rm lab}\|_{C^2}\le Ch^{s-2}.
\end{equation}
Here uncentered support functions are compared in the same laboratory
frame; centering preserves these orders.
\end{lemma}
\begin{proof}
Use the support function $p$ relative to $o$, so $p\ge r$. In normal
angle $\theta$, the boundary is $p n+p' n^\perp$, its radial angle is
\[
 \varphi=\theta+\arctan(p'/p),\qquad
 \frac{\dd\varphi}{\dd\theta}
             =\frac{p(p+p'')}{p^2+(p')^2}>c>0.
\]
The inverse $\theta(\varphi)$ is uniformly $C^{5,\beta}$. Although
the expression $R=\sqrt{p^2+(p')^2}$ initially gives only five
ordinary derivatives, differentiation in radial angle cancels the
curvature factor:
\begin{equation}\label{eq:radial-cancellation}
                   R'=R(p'/p)\circ\theta.
\end{equation}
Its right side is uniformly $C^{5,\beta}$. Thus $R$ is uniformly
$C^{6,\beta}$, without losing a derivative from the support prior.

At each angular node interpolate the seven consecutive sampled
values by a polynomial of degree at most six in a local lift of the
angle. Combine these polynomials with a fixed translated smooth
partition of unity supported on neighborhoods of size $O(h)$.
On each such neighborhood the normalized Vandermonde inverse is
bounded. A degree-six Taylor polynomial of the true function is
reproduced exactly, and its remainder and derivatives through order
three are $O(h^{s-j})$. Errors in the values contribute $O(eh^{-j})$.
Derivatives falling on the partition have size $O(h^{-j})$ and are
absorbed by the corresponding lower-order remainders. The overlap
number is fixed. This proves \eqref{eq:radial-interpolation}, including
across the angular cut by periodic lifts.

A positive radial function is strictly convex when
\[
       F=R^2+2(R')^2-RR''>0.
\]
For the true curve this expression has a uniform positive margin,
by the inradius, curvature and diameter bounds. It persists for
$\widehat R$ by the $C^2$ estimate. The two radial graphs are uniformly
$O(h^s)$-close, which proves the Hausdorff bound.

To check the stronger support norm, set
$\theta=\varphi-\arctan(R'/R)$, now viewing the normal angle as a
function of radial angle. Its derivative is
$F/(R^2+(R')^2)>c$. At corresponding angles the support value and
radius of curvature are
\[
 p=\frac{R^2}{\sqrt{R^2+(R')^2}},\qquad
 p+p''=\frac{(R^2+(R')^2)^{3/2}}{F}.
\]
These expressions, and the tangential support derivative, depend
Lipschitzly on $R,R',R''$ on the bounded set with these margins.
The inverse normal-angle maps differ by $O(\|\widehat R-R\|_{C^1})$.
Their composition changes the true expressions by the same order,
using the uniform $C^3$ bound; the reconstructed functions also have
such a bound by \eqref{eq:radial-interpolation}. Thus the support
error through two derivatives is $O(h^{s-2})$. Translation to the
laboratory frame adds the same first harmonic to both supports. The
Steiner point is a bounded linear functional of the support, so its
error is $O(h^s)$ and centering preserves the conclusions.
\end{proof}

\begin{proposition}[Adaptive patch reconstruction]\label{prop:patch}
For a fixed protected patch, complete component boundaries can be
recovered with the errors in \eqref{eq:geometric-interpolation} using
at most $C h^{-1}\log(C/h)$ local occupation queries. The finest
localization is $\sigma=c h^s$. With probability at least $1-\delta$
the attempted-bit cost is at most
\begin{equation}\label{eq:patch-cost}
           C h^{-1}\log(C/h)\log(C/(h\delta)).
\end{equation}
The number of components, the initial coarse queries and every
local Bellman stencil size are bounded solely by the fixed prior.
\end{proposition}
\begin{proof}
On the simultaneous accuracy event, obtain the hulls and centers in
Lemma~\ref{lem:coarse}. Separation and the fixed aperture bound their
number. For each body use $m$ equally spaced radial directions and
their isolated coarse brackets. Lemma~\ref{lem:bisection} uses
$O(\log(C/h))$ midpoint queries per direction to return values with
error $e\le C\sigma$. Lemma~\ref{lem:radial} gives the claimed
geometric estimates. The total number of queries is bounded by a
deterministic $Q_h=C(1+h^{-1}\log(C/h))$. Allocate failure probability
$\delta/Q_h$ to every query, including the coarse ones, before the
experiment. Conditional confidence and the union bound in
Proposition~\ref{prop:query} give the simultaneous event. Its
preparation cost is \eqref{eq:patch-cost}. Locations chosen using
coarse or previous fine data do not alter this argument.

For a finite numerical realization, polygon intersections, midpoint
locations and radial values are enclosed with error $O(\sigma)$.
The estimate \eqref{eq:radial-interpolation} absorbs those value errors.
Represent the output by its finite polynomial coefficients and the
fixed smooth partition functions. They, their derivatives and the
normal-angle maps can be evaluated to any positive prescribed error
by finite operations on compact intervals with the stated derivative
bounds. This constructs an actual smooth convex output, without
assigning a false smoothness bound to the initial polygons. No
infinite-dimensional candidate minimization is performed.
\end{proof}

\section{Finite binary descriptions of the adaptive controls}
\label{sec:digital}
The active protocol specifies a distribution by a finite command, rather
than transmitting an exact real coordinate. We separate this digital
resource from the physical production of a localized launch distribution.
The latter remains a calibrated apparatus primitive; it is not obtained
by observing the collision bits.

\begin{theorem}[A dyadic controller and its resource bounds]
\label{thm:digital}
On the same bounded class $\mathcal B_\eta$, choose the compass length
$t>0$ and the fixed coarse numerical reserves dyadic, while retaining
all strict separation margins. The experiment of
Theorem~\ref{thm:adaptive} can be implemented with dyadic nominal
centers and dyadic localization scales. Its accuracy, deterministic
attempt bound and exact discrete outputs are unchanged up to constants.
Put
\[
 b=\lceil C_0+\log_2(1/\sigma)\rceil,
 \qquad B=1+\lceil\log_2(C/\nu)\rceil
              +\lceil\log_2(C/\delta)\rceil .
\]
The finest center mesh is $2^{-b}\le c\sigma$, with
$b=O(\log(C/\nu))$. Each batch command has a binary description of
length at most $CB$, and the total transmitted command description is
at most $CJ_\nu B$. The finite boundary representation has length
at most $C\nu^{-1/(s-2)}B$. These bounds include batch repetition
counts. They do not count the analog random numbers internal to the
prescribed launch distribution as table-dependent observations.

For rational fixed numerical priors and rational $\beta$, one may
choose the fixed smooth partition and its evaluators so that a
conservative digital bit-operation bound is
\begin{equation}\label{eq:digital-work}
              C(N_\nu+J_\nu+\nu^{-1/2})B^c
\end{equation}
for an absolute finite exponent $c$. This includes boundary output,
primitive arithmetic and free-area evaluation to error $C\nu$.
It is not a sharp computational complexity bound. For general fixed
$\beta$ the message bounds hold with certified dyadic scale enclosures;
a bit-operation bound additionally depends on their computation cost.
\end{theorem}
\begin{proof}
\emph{One exact dyadic stencil.}
Take the fine localization to be a dyadic number comparable to $c h^s$,
where $h=2\pi/m\asymp\nu^{1/(s-2)}$. Choose a dyadic rectangle $W$
with all the protective collars of Lemma~\ref{lem:stopped}. Round a
requested target $x$ once to a dyadic center $x'$, with
$|x-x'|\le C2^{-b}$. Use the exact stencil
\[
                         x'+t(i,j),\qquad |i|+|j|\le n_*.
\]
For $b$ greater than the fixed denominator exponent of $t$, all stencil
points are on the same dyadic grid. Membership in $W$ is tested by
exact rational comparisons. There is no independent rounding of
transition targets, hence no error in the killed compass identity.
The reverse nominal distribution is exactly the translated joint law:
it attempts $(q+a,-a)$ for the same specified law of $(q,a)$, not an
independently rounded reverse start. Fresh attempts are still used.

The query estimates $u_\sigma(x')$. Its derived label at the intended
point $x$ is therefore correct outside the boundary tube of width
$\sigma+C2^{-b}$. Physical position, duration and angle errors receive
an independent portion of the reserve in
Proposition~\ref{prop:query}. Thus rounding and actual calibration
together preserve its fixed accuracy bound. The means are ratios of
integer hit counts to a deterministic batch size. Taking that size to
be a common integer upper bound for all queries makes every local
average and Bellman update an exact rational computation. Clipping and
comparison with $1/2$ are also exact. The depth and stencil size are
prior constants.

\emph{Rounded adaptive geometry.}
The coarse grid and its finite polygons have a prior-bounded number
of rational vertices. A dyadic interior center with a reserved fixed
inball can be found by finite search; squared distances to rational
sides test the inball inequalities. The strict inball reserve ensures
that a sufficiently fine fixed search grid succeeds. Approximate
$\cos(2\pi j/m)$ and $\sin(2\pi j/m)$ with certified error at most
$c2^{-b}$. Polygonal radial values and brackets are computed with the
same reserve. Consequently every implemented spatial midpoint is
within $C2^{-b}$ of the intended radial point.

For completeness, rounding a radial midpoint by at most $d$ does not
accumulate $kd$ after $k$ queries. The relaxed interval invariant from
Lemma~\ref{lem:bisection} holds with a boundary layer of width
$A(\sigma+d)$. Its interval widths satisfy
\[
                     w_{k+1}\le w_k/2+d,
 \qquad w_k\le2^{-k}w_0+2d.
\]
Stop before asking for resolution finer than the mesh. With
$d=C2^{-b}\le c\sigma$, $O(\log(C/\sigma))$ queries give radial
values of error $O(\sigma)$, even when labels in the enlarged layer
are chosen adaptively. Round the final values to the same mesh.
Lemma~\ref{lem:radial} absorbs these errors and still gives actual
smooth strictly convex output bodies. The known patch margin and
bounded-denominator gaps allow the same exact period decisions.

\emph{Description lengths.}
All centers lie in a fixed rectangle. Their dyadic integer numerators
therefore use $O(b)$ bits. A batch command consists of a forward/reverse
flag, two such coordinates, a scale exponent, the fixed compass/kernel
identifier, and a repetition count of order $\log(C/(\nu\delta))$.
The last count requires only $O(\log B)$ binary digits. Fixed-width
fields give the asserted $CB$ bound without a real-valued control
field. There are at most $CJ_\nu$ batches, allowing repetitions.
The receiver interprets a fixed kernel identifier as the stipulated
localized random preparation, not as a table-dependent data stream.

Store the dyadic radial values, the dyadic centers and the integers
$j,m$. In normalized coordinates the seven-point interpolation matrix
has fixed rational entries, so its coefficients have $O(B)$-bit
rational descriptions. Local angles and partitions are generated by
a fixed program from $j,m$; no table-dependent exact trigonometric
constant is stored. There are $O(m)$ coefficient blocks, proving the
boundary-description bound.

\emph{A conservative operation count.}
For the computational assertion choose an explicit smooth bump,
for example $\exp[-1/(1-x^2)]$ on $|x|<1$, zero outside, and normalize
a finite-overlap family of its translates to make the fixed partition.
The denominator has a positive fixed lower bound. To the required
precision, the bump and its first three derivatives are evaluated by
rational series on bounded arguments; near its endpoints an explicit
exponential tail bound permits certified truncation to zero. The
trigonometric functions, $\pi$, square roots and logarithms needed
here admit rational series and bisection with a polynomial number of
bit operations in $B$. For rational $s$, dyadic powers comparable to
$h^s$ can equivalently be bracketed by raising rational quantities to
fixed integer powers. Thus all scale choices are effective in this
computational subcase.

Processing bit counts, the constant-depth rational Bellman updates,
rounded bisection and the fixed-degree coefficient blocks costs
$C(N_\nu+J_\nu)B^c$. One need not optimize global numerical work to
obtain a bound: the reconstructed support functions have uniform
$C^2$ bounds. Composite trapezoidal quadrature with
$O(\nu^{-1/2})$ nodes computes their Steiner integrals, and the area
integrals $\tfrac12\int\widehat R^2$, to error $O(\nu)$. A support
value is evaluated by inverting the uniformly monotone normal-angle
map; bisection with residual enclosures uses $O(B)$ steps. Each
partition evaluation involves only a fixed number of neighboring
coefficient blocks. The uniform derivative and convexity margins
justify all these tolerances.

The number of protected components and candidate period pairs is a
prior constant. Support-defect comparisons need only a fixed precision
smaller than the known patch and rational-separation margins, whereas
the real period coordinates and areas are evaluated to $O(\nu)$.
Their mesh tests, bounded-denominator rational search and two-dimensional
Hermite arithmetic are finite with polynomial cost in $B$. Including
the preceding quadratures proves \eqref{eq:digital-work}. The larger
$\nu^{-1/2}$ term is a disclosed conservative numerical overhead; it
is not silently identified with the attempted-bit exponent.
\end{proof}

\subsection{Separate resource axes}
Write $m_\nu=\nu^{-1/(s-2)}$ and $L_\nu=\log(C/\nu)$. The following
accounting refers to the construction, not to a measured apparatus.
\begin{center}
\small
\begin{tabular}{@{}p{0.31\textwidth}p{0.64\textwidth}@{}}
\toprule
Resource & Bound or explicit status\\
\midrule
Table-dependent outputs & $N_\nu\le C m_\nu L_\nu\log(C/(\nu\delta))$ bits;
all attempts, including zeros, are charged.\\[3pt]
Command centers/batches & $J_\nu\le C m_\nu L_\nu$, counted with repetitions.\\[3pt]
Digital controls & $O(B)$ bits per batch; $O(J_\nu B)$ in total.\\[3pt]
Spatial localization & $\sigma\asymp\nu^{s/(s-2)}$; the localized kernel
is a prescribed preparation primitive.\\[3pt]
Physical calibration & $\ell+\tau+t\alpha\le c\sigma$ is sufficient.
Section~\ref{sec:calibration-floor} gives a matching necessary position-error power.\\[3pt]
Aperture & Fixed by the geometric prior, including reverse starts and
protective collars.\\[3pt]
Numerical bit operations & The conservative bound \eqref{eq:digital-work}
under its explicit effective-prior assumptions.\\[3pt]
Motion, setup, calibration production & No physical time, energy or
cost law is assumed. No bound on these resources is inferred from $N_\nu$.\\
\bottomrule
\end{tabular}
\end{center}
Digital description length is not mechanical resolution. In particular,
a finite word specifying a narrow kernel does not prove that a device
can produce that kernel cheaply. The converse calibration theorem
below quantifies a necessary tolerance in the observation model,
without inventing a law for the cost of achieving it.

\section{Period recognition from the reconstructed patch}
\label{sec:periods}
This section retains the protected-patch argument for repeated shapes.
It acts on the acquired geometry, not on an assumed list of integer
copy labels. We include the proof to specify exactly which discrete
information the adaptive experiment recovers.

Put $B=1+L_0$ and $Q_a=[-a,a]^2$. The aperture used above is chosen
large enough to recover complete components with centers in $Q_{10B}$,
together with their protective collars. For two components define
\[
 d_*(C,D)=\|z_C-z_D\|_\infty+\|p_C-p_D\|_\infty,
\]
where the centered support functions are compared in the laboratory
directions, without an independent rotation. Define the patch defect
\begin{equation}\label{eq:defect}
 \mathcal D(w)=\max_{C:z_C\in Q_B}\ \max_{\epsilon\in\{-1,1\}}
                     \inf_D d_*(C+\epsilon w,D).
\end{equation}
The infimum ranges over physical components. Distant components cannot
minimize it. The known nonperiod margin is the requirement
\begin{equation}\label{eq:patch-margin}
 \begin{gathered}
 w=z_D-z_C,\quad z_C\in Q_{3B},\ z_D\in Q_{5B},\quad w\notin\Pi
       \quad\Longrightarrow\quad\mathcal D(w)\ge\eta>0.
 \end{gathered}
\end{equation}
This is a margin for whole-patch translations, not a separation between
all individual shapes. Every fixed table in $\mathcal B$ has some
positive such margin, since there are finitely many candidate pairs.
The margin is not claimed to be uniformly inferable from finite bits.

\begin{lemma}[Exact patch recognition]\label{lem:recognition}
One has $\mathcal D(w)=0$ if and only if $w\in\Pi$. For
$\|w\|_\infty\le8B$, this zero test uses only components centered in
$Q_B$ and $Q_{9B}$. The group $\Pi$ is a lattice with
\begin{equation}\label{eq:index-bound}
 [\Pi:L\Z^2]\le r_0,\qquad\covol\Pi\ge V_*=V_0/r_0.
\end{equation}
Choose any component centered in $Q_{2B}$ as a root. Its center
differences to components centered in $Q_{4B}$ that pass the zero
test generate $\Pi$ over $\Z$.
\end{lemma}
\begin{proof}
Each orbit modulo $L\Z^2$ has a representative centered in
$L[-1/2,1/2]^2\subset Q_B$. A zero defect gives both translates of
every such representative as physical components. Translating by
$L\Z^2$ gives $\mathcal O+w\subset\mathcal O$ and
$\mathcal O-w\subset\mathcal O$; together these imply equality.
The converse is immediate. The centers of a matching pair in the
bounded test are in $Q_{9B}$.

Map a coset of $\Pi/L\Z^2$ to the component orbit of the translated
root. This is injective: a nonempty compact body is not invariant
under a nonzero translation. Thus the quotient has at most $r$ elements.
A finite extension of a full-rank lattice is a lattice, giving
\eqref{eq:index-bound}. The two columns of $L$, and representatives
of all these period cosets, can be chosen with norm less than $B$
in the maximum norm. Translating the root by these vectors gives
targets in $Q_{3B}$, strictly within the cutoff $Q_{4B}$. They are
included in the accepted list and generate $\Pi$. Conversely each
accepted vector is a true period.
\end{proof}

\begin{lemma}[Stable discrete decisions]\label{lem:locking}
Suppose the protected complete components have matched Hausdorff
error at most $e$. On $\mathcal B_\eta$, for sufficiently small $e$
depending on the full prior, finite tests recover the true primitive
orbit count and exact rational relations of a true accepted generating
list in a true independent-pair basis. They recover a matched period
basis and primitive free area to error $O(e)$.
\end{lemma}
\begin{proof}
Hausdorff error $e$ gives Steiner-center error at most $2e$ and
centered-support error at most $3e$. Use approximate source centers
in $Q_{2B}$, a root there and root targets in $Q_{4B}$. For a measured
difference $\widehat w$, its corresponding true difference has error
at most $4e$. Test both signs on all selected sources, minimizing over
targets with approximate centers in $Q_{10B}$. The center contribution
to each comparison changes by at most $8e$ and the support contribution
by at most $6e$. True periods therefore have measured defect at most
$14e$. For a nonperiod, its true endpoints lie in $Q_{3B},Q_{5B}$,
so \eqref{eq:patch-margin} supplies a witness of defect at least $\eta$.
All witnesses centered in $Q_B$ are selected. Restricting targets cannot
lower the true infimum. Its measured defect is at least $\eta-14e$.
Reserve another $6e$ for numerical enclosures and use threshold $\eta/2$,
with $20e<\eta/4$. This accepts exactly the true periods. Strict cutoff
collars ensure that every generator in Lemma~\ref{lem:recognition}
is present. The same tests on pairs of sources give their primitive
orbit relation: a period translating the first center to the second
must translate the first body to the second body, since Steiner points
lie in the interiors and different components are disjoint. All
primitive orbits are represented, so their number is recovered.

Accepted true vectors have a prior bound $U_0$ on their Euclidean
lengths. A nonzero determinant of two such periods is an integer
multiple of $\covol\Pi\ge V_*$. Determinant perturbations are
$O(U_0e)$; hence independence is decided below this gap. Select an
independent true pair with matrix $D$. Its generated sublattice has
index
\[
       n=|\det D|/\covol\Pi\le Q:=\max(1,\lceil U_0^2/V_*\rceil).
\]
Every coordinate of $D^{-1}w_j$ has reduced denominator dividing $n$
and a prior-bounded numerator. The inverse-matrix identity bounds its
estimation error by $C_*e$. Distinct rationals of denominators at most
$Q$ are separated by at least $Q^{-2}$. Finite search locks the exact
coordinates when $C_*e<1/(3Q^2)$.

Let $q$ be their common denominator, which divides $n$. A column
Hermite basis $H_0$ of the integer span of $qI_2$ and $qD^{-1}w_j$
gives
\begin{equation}\label{eq:period-basis}
                         \Pi=(DH_0/q)\Z^2.
\end{equation}
The integer group contains $q\Z^2$, so its Hermite entries have a
prior bound. Replacing $D$ by its estimate gives a matched true basis
with error $O(e)$. We have recovered exact integer relations, not
exact real coordinates. The primitive free area is
\[
             A_* =\covol\Pi-\sum_{j=1}^{r_*}\area(C_j).
\]
Uniform perimeter bounds and the parallel-body inclusions give area
error $O(e)$ for each matched component, proving the last assertion.
Support suprema and center integrals used in the tests admit certified
finite angular meshes, because the bounded bodies have a uniform
support Lipschitz bound. Their numerical errors are charged above.
\end{proof}

\begin{proof}[Proof of Theorem~\ref{thm:adaptive}]
Choose a sufficiently large fixed target square so that every component
needed in Lemma~\ref{lem:locking} and every coarse bracket is protected.
Separation bounds their number. Choose the fixed observation and launch
apertures as in Lemma~\ref{lem:stopped}. Apply
Proposition~\ref{prop:patch} with an integer angular mesh satisfying
$h\asymp\nu^{1/(s-2)}$ and a sufficiently small constant factor.
Its reconstructed bodies have matched $C^2$ support error $O(\nu)$
and Hausdorff error $e=O(h^s)$. For small $\nu$, all thresholds in
Lemma~\ref{lem:locking} hold. It gives exact primitive discrete data
and period-basis error $O(h^s)$. Repeating one recovered body per
primitive orbit under this basis gives a physical periodic presentation.
Indeed the inverse basis remains bounded, so only a bounded number of
integer translates can enter a bounded fundamental patch. The positive
separation and curvature margins then persist. The reconstructed
bodies have area error $O(h^s)$ and the covolume has the same order;
hence the primitive free-area error is also $O(h^s)$, in particular
$O(\nu)$.

The query count of Proposition~\ref{prop:patch} and its fixed local
stencil size give \eqref{eq:centers}. Equation~\eqref{eq:patch-cost}
gives \eqref{eq:attempts}, since $s$ is fixed. The calibration bias
allocation in Proposition~\ref{prop:query} gives \eqref{eq:precision}.
Every preparation has already been charged there. All steps are finite,
and the output is a finite smooth representation with certified
numerical tolerances. This proves the theorem.
\end{proof}

Without a known $\eta$, successive refinements give the corresponding
pointwise eventual decisions for each fixed table, but not a uniformly
valid observable stopping certificate. The exact functional inverse
combined with Lemma~\ref{lem:recognition} needs no supplied positive
margin: it tests zeros exactly. This distinction, as well as the
bounded periodic presentation itself, is retained. No inference of
crystallinity from an arbitrary nonperiodic configuration is made.

\section{A physical packing and a binary-transcript bound}
\label{sec:lower}
The lower bound counts the table-dependent binary outputs, not the
precision of the controller's choices. It therefore applies to more
powerful binary command designs as well as the reciprocal protocol.
Its geometric construction has fixed periods and a fixed positive
patch margin, so period uncertainty is not the cause of the bound.

\begin{lemma}[A uniform physical packing]\label{lem:packing}
For $s=6+\beta$ there is a fixed class $\mathcal P\subset\mathcal B_\eta$
and constants $c_0,c_1>0$ such that, for all sufficiently small $h$,
$\mathcal P$ contains $2^m$ tables, where $m\ge c_0/h$, whose variable
component support functions have pairwise $C^2$ distance at least
$c_1h^{s-2}$. All have the same known primitive lattice and two
primitive component orbits. Their correspondences and pose can be
supplied without reducing this separation.
\end{lemma}
\begin{proof}
Fix the lattice $20\Z^2$. Put one variable component near the unit
disk at the origin, and a fixed disk of radius two at $(6,0)$, and
repeat both by this lattice. Choose a nonzero smooth function
$\psi$ supported in $(-1/3,1/3)$ with $\psi''(0)\ne0$.
Take angular centers $\theta_j$ in a fixed interval with mutual
spacing at least $h$, with $c/h\le m\le C/h$. For
$\omega\in\{0,1\}^m$, define the uncentered support of the variable
body by
\begin{equation}\label{eq:packing}
 p_\omega(\theta)=1+a h^s
                   \sum_{j=1}^m\omega_j\psi((\theta-\theta_j)/h),
\end{equation}
periodically extended, where $a>0$ is a fixed sufficiently small
number. At most one summand is nonzero at each point. Its derivative
of order $k\le6$ is bounded by $Ca h^{s-k}$. The sixth derivative
has uniformly bounded $\beta$-H\"older seminorm. Within a support
this follows by rescaling; between supports it follows from their
separation and the vanishing at their edges. The same estimate follows
by comparison to the intervening zero point for two arbitrarily close
points on different sides of an edge. Thus the $C^{6,\beta}$ bound
is uniform over all $h,m,\omega$.

For small $a$, $p_\omega+p_\omega''$ lies in a fixed positive interval.
The bodies are smooth and strictly convex, with all separation,
diameter, lattice and curvature bounds independent of $h$. Subtracting
their Steiner first harmonic also preserves the uniform support prior.
The two component species cannot be translated into each other:
their diameters lie in disjoint fixed intervals. Every period must
therefore send a variable component to another variable component,
forcing that period to lie in $20\Z^2$. Conversely every such vector
is a period. This proves primitiveness and the stated orbit count.

Center differences of the same species are true periods. For any
cross-species difference, the translated source species cannot match
a target of the other species with small support error, because their
widths differ by a fixed amount. A target of the same species has
center discrepancy bounded below by a fixed positive number: the
cross-species offset $(6,0)$, perturbed by the variable Steiner point,
stays a positive distance from $20\Z^2$. There are representatives
of both species in $Q_B$. These two alternatives give a uniform
positive defect for every nonperiod candidate in
\eqref{eq:patch-margin}. Hence one fixed $\eta>0$ works throughout
the construction. Taking the union of these tables over all small
$h$ defines a single class $\mathcal P$ with fixed numerical bounds.

If $\omega\ne\omega'$, choose an index where they differ and evaluate
the second derivative at its center. All other summands vanish there,
so the absolute difference is $a|\psi''(0)|h^{s-2}$. This proves the
separation in the laboratory support norm. The variable component,
lattice and fixed component may be identified in advance; none of
these disclosures identifies its bump vector. No analytic smoothness
bound is used or claimed for this packing.
\end{proof}

\begin{lemma}[Binary decision trees with randomized commands]
\label{lem:bits}
Let $M$ parameters have pairwise disjoint acceptable output sets.
A procedure with at most $N$ binary table-dependent outputs, arbitrary
adaptive inputs and parameter-independent randomization has average
success probability at most $2^N/M$ under the uniform distribution on
these parameters. The binary response probabilities may be arbitrary
functions of the parameter and chosen input.
\end{lemma}
\begin{proof}
Represent the controller randomization by an independent seed. Each
binary response can be generated by comparing another independent
uniform random variable with its conditional response probability.
This constructs the joint adaptive experiment for every parameter on
one common seed space. Fix all these independent seeds. Commands are
then functions only of the preceding binary transcript, and the output
is a function of a binary tree with at most $2^N$ leaves, after padding
short transcripts if necessary. An output at a leaf is acceptable for
at most one parameter, by disjointness. Thus at most $2^N$ parameters
can succeed for these fixed seeds. Average over the uniform parameter
and then integrate over the seeds. The bound follows. The numerical
values of selected commands give no additional leaves, since they are
already determined by the seed and transcript.
\end{proof}

\begin{proof}[Proof of Theorem~\ref{thm:lower}]
In Lemma~\ref{lem:packing} take
$h\asymp\nu^{1/(s-2)}$ with its constant large enough that the support
separation exceeds $2\nu$. The $\nu$-accuracy output sets are disjoint.
Lemma~\ref{lem:bits}, with $M=2^m$, gives average success at most
$2^{N-m}$. Uniform success at least $3/4$ therefore requires
$N\ge m+\log_2(3/4)\ge c\nu^{-1/(s-2)}$ for small $\nu$.
This argument allows more input control than the upper-bound procedure
and assumes perfect calibration. It remains valid for the latter
procedure with all-attempt bookkeeping. A protocol that reveals a
free-start test or a continuous impact location has a different output
alphabet and is not covered by the assertion.
\end{proof}

The loss in this theorem is the supplied laboratory-frame loss used by
the active controller. It is already enough to establish a matching
power for Theorem~\ref{thm:adaptive}, in the same loss. We do not silently
replace it by an unlabelled quotient loss in the lower bound, nor extend
a deterministic cap $N$ to an expected stopping-time budget without
another argument. No lower bound for the logarithms or for confidence
dependence is asserted.

For fixed confidence, on a nondegenerate bounded class containing
$\mathcal P$, the two results give
\begin{equation}\label{eq:matched-power}
 c\nu^{-1/(4+\beta)}\le N^*(\nu)
          \le C\nu^{-1/(4+\beta)}\log^2(C/\nu),
\end{equation}
where $N^*$ is the least uniform deterministic cap on attempted bits,
and constants in the accuracy definition are absorbed by rescaling
$\nu$. The matching power comes from one-dimensional boundary
smoothness measured two derivatives below the prior. The physical
work needed for the upper bound is that these boundary queries can
be implemented from the stipulated collision bits without additional
geometric outputs or uncharged free preparations.

\section{The lower-bound gauge and sequential observations}
\label{sec:sequential}
The deterministic-cap result above is retained. We now make its gauge
explicit and prove a separate result for expected stopping times. No
real-valued stopping signal or unobserved free-start test is allowed.

\subsection{A parameter-independent laboratory frame}
Let $n(\theta)=(\cos\theta,\sin\theta)$. For a laboratory support $p$
write
\begin{equation}\label{eq:centering-map}
 z(p)=\frac1\pi\int_0^{2\pi}p(\theta)n(\theta)\dd\theta,
 \qquad p^0=p-z(p)\cdot n.
\end{equation}
The disclosed information in Lemma~\ref{lem:packing} consists of the
fixed laboratory axes, $20\Z^2$, the names of the two component species,
and the fixed radius-two disk at $(6,0)$. These data are identical for
all packing parameters. In particular, the unknown Steiner point of
the variable body is not disclosed. ``Pose given'' does not mean a
parameter-dependent translation or rotation of that body.

For the variable component define acceptable output sets by
\begin{equation}\label{eq:accuracy-sets}
 \mathcal U_\omega(\nu)=
 \{\widehat C:\ \|p^0_{\widehat C}-p^0_\omega\|_{C^2}\le\nu\}.
\end{equation}
This centered-shape loss alone is enough for the lower bound. The
stronger laboratory loss may also include
$|z_{\widehat C}-z(p_\omega)|$ and errors of the fixed component.
No independent rotation of each output is taken in either loss.

\begin{lemma}[The packing survives centering]\label{lem:centered-packing}
For the family \eqref{eq:packing}, distinct parameters satisfy
\[
 \|p^0_\omega-p^0_{\omega'}\|_{C^2}\ge c h^{s-2}
\]
for all sufficiently small $h$, uniformly over the packing. The
centered physical support bounds and the patch margin remain uniform.
\end{lemma}
\begin{proof}
The disjoint angular supports give
$\|p_\omega-p_{\omega'}\|_\infty\le C a h^s$, independently of
$m$. Equation~\eqref{eq:centering-map} implies
$|z(p_\omega)-z(p_{\omega'})|\le2Ca h^s$. The $C^2$ norm of the
subtracted first harmonic is therefore $O(a h^s)$. At a differing
bump center the uncentered second-derivative difference has magnitude
$a|\psi''(0)|h^{s-2}$. Since $h^s=o(h^{s-2})$, at least half this
difference survives centering. A first harmonic is bounded in every
fixed smoothness norm by its amplitude, so the centered prior is also
preserved. The separation and period arguments of
Lemma~\ref{lem:packing} already allow these $O(a h^s)$ center shifts.
Thus the disclosed data contain no shape-dependent centering information.
\end{proof}

\subsection{A stopped transcript carries at most its expected length}
A stopping rule is admissible when, conditional on parameter-independent
controller randomness, stopping, further commands and the eventual
output depend only on the preceding bits. Waiting times or a physical
success flag are not additional data. The number $\mathsf T$ counts
all attempted preparations, including zeros from solid starts.

\begin{lemma}[Sequential binary information bound]\label{lem:sequential}
Let $M\ge2$ parameters have disjoint acceptable output sets. If an
admissible randomized procedure succeeds with probability at least
$1-\delta$ for each parameter, $0<\delta<1/2$, then
\begin{equation}\label{eq:sequential-information}
 \frac1M\sum_{j=1}^M\E_j\mathsf T
 \ge (1-\delta)\log_2M-h_2(\delta),
\end{equation}
where $h_2(x)=-x\log_2x-(1-x)\log_2(1-x)$. There is no deterministic
cap in this assertion. An infinite expected length satisfies it
trivially.
\end{lemma}
\begin{proof}
Give the parameter $V$ the uniform distribution and denote the
independent controller seed by $U$. Suppose the average expected length
is finite, so the terminal binary word $Z$ is finite almost surely.
For each fixed $U=u$, terminal words are prefix-free: after the rule
stops at one word it cannot also stop at a proper extension. Thus
$\sum_z2^{-|z|}\le1$. If this sum is $K_u>0$, comparison with the
probability distribution $2^{-|z|}/K_u$ by nonnegativity of relative
entropy gives
\[
 H(Z\mid U=u)\le\E(|Z|\mid U=u)+\log_2K_u
                    \le\E(\mathsf T\mid U=u).
\]
The countable case follows from the same log-sum inequality by
truncation. Integrate in $u$. Since $U$ is independent of $V$,
\[
 I(V;Z,U)=I(V;Z\mid U)\le H(Z\mid U)\le\E\mathsf T.
\]
Decode the parameter from the unique acceptable output set, using an
arbitrary index on failure. Its average error $p_e$ is at most $\delta$.
For the error indicator $E$, the elementary conditional-entropy bound is
\[
 H(V\mid Z,U)\le H(E\mid Z,U)+p_e\log_2(M-1)
 \le h_2(\delta)+\delta\log_2M.
\]
Subtract from $H(V)=\log_2M$ and use the preceding information bound.
The output randomness may be included in $U$ at the outset. Neither
continuous input values nor the word length yield an extra channel:
they are already determined by the seed and stopped word. The coding
and entropy principles used here are classical \cite{Shannon}; the
argument is included to specify the stopping-time scope exactly.
\end{proof}

\begin{theorem}[Expected attempted-bit lower bound]\label{thm:expected}
On the fixed physical class of Lemma~\ref{lem:packing}, any admissible
sequential procedure with success probability at least $3/4$ uniformly
in the centered loss \eqref{eq:accuracy-sets} satisfies
\begin{equation}\label{eq:expected-lower}
                 \sup_{\mathcal O}\E_{\mathcal O}\mathsf T
                         \ge c\nu^{-1/(s-2)}.
\end{equation}
The bound holds with exact implementation and arbitrary-precision
adaptive inputs, even with the fixed lattice, species and laboratory
frame disclosed. It therefore also holds for the stronger physical
reconstruction loss.
\end{theorem}
\begin{proof}
Choose $h\asymp\nu^{1/(s-2)}$ so that Lemma~\ref{lem:centered-packing}
gives separation greater than $2\nu$. Apply Lemma~\ref{lem:sequential}
with $M=2^m$ and $m\ge c/h$. The supremum dominates the average.
\end{proof}

At fixed confidence, the least uniform expected number of attempts on
any fixed class containing this packing lies between
$c\nu^{-1/(s-2)}$ and
$C\nu^{-1/(s-2)}\log^2(C/\nu)$, since the deterministic upper bound
is also an expected upper bound. The analogous deterministic-cap
statement remains \eqref{eq:matched-power}. These are two explicitly
proved resource criteria, not an identification of random length with
a deterministic cap. No sharp logarithmic or confidence factor follows.

\section{A resolution converse for table-dependent bounded errors}
\label{sec:calibration-floor}
The sufficient calibration scale in \eqref{eq:precision} has a converse
in the same bounded-error model. The converse uses physical start
perturbations, not corrupted labels added to an abstract membership
oracle. It concerns short collision commands; arbitrarily long flights
or additional observations of the actual launch position are not included.

\subsection{A common response produced by nearby physical tables}
For a table $\mathcal O_j$ denote its attempted hit bit by $B^j_a(q)$.
An admissible position-error implementation is a measurable map
$q\mapsto F_j(q,a)$ with $|F_j(q,a)-q|\le r$; time and direction
can remain exact. Such maps may depend on the table and the nominal
command. The controller sees only the resulting bit, not the actual
start. This is a submodel of the conditionally bounded, not necessarily
mean-zero couplings in Section~\ref{sec:setting}.

\begin{lemma}[Pointwise common response under two implementations]\label{lem:common-response}
Suppose $\mathcal O_0,\mathcal O_1$ consist of matched strictly convex
components of diameter at most $D_0$, with inradius at least $r_*>0$,
and that matched bodies have Hausdorff distance at most $e>0$.
In each table distinct bodies are separated by at least $d_0$.
Fix $t_{\max}<d_0$. If $e$ is sufficiently small compared with
$r_*$ and $d_0-t_{\max}$, there are measurable position perturbations
of size at most $2e$ such that, for every $|a|\le t_{\max}$ and
with closed-solid and swept-set conventions,
\begin{equation}\label{eq:common-response}
 B^0_a(F_0(q,a))=B^1_a(F_1(q,a))
                    =\min\{B^0_a(q),B^1_a(q)\}.
\end{equation}
The closed convention assigns zero on a solid boundary and one on a
swept boundary outside the solid. In every disagreement case the
forced-zero position is moved strictly into a solid or strictly outside
all sweeps before its response is evaluated. Agreement cases keep the
nominal point and use the same closed convention. The conclusion
applies to locally finite periodic arrays and to both the forward and
translated reverse commands.
\end{lemma}
\begin{proof}
For one fixed $a$ put $S_j(C)=C+[-a,0]$. Matched swept bodies have
Hausdorff distance at most $e$, because their support functions differ
by the same amount as those of the matched bodies. In either table,
\begin{equation}\label{eq:swept-separation}
 \dist(S_j(C),S_j(D))\ge d_0-|a|\quad(C\ne D).
\end{equation}
Indeed two points in the sweeps are $c-ua$ and $d-va$, with
$u,v\in[0,1]$. Their distance is at least $|c-d|-|u-v||a|$.
Thus a start belonging to a swept body determines that body uniquely.
Moreover no solid other than that body meets this sweep. On it the
hit bit is the indicator of the sweep minus the indicator of its body.

Leave the start unchanged when the two bits agree, and also on the
side whose bit is already zero. Consider $B^0_a(q)=1$, $B^1_a(q)=0$,
and let $C_0$ be the unique body whose sweep contains $q$. Write $C_1$
for its matched body. The separation \eqref{eq:swept-separation} and
Hausdorff matching imply the following exhaustive alternatives:
\[
                   q\in C_1,\qquad\hbox{or}\qquad q\notin\operatorname{int}S_1(C_1).
\]
To check this, every other sweep in table 1 is at distance at least
$d_0-|a|-e$ from $S_0(C_0)$, so cannot supply either a competing hit
or a solid start there. If $q\in S_1(C_1)\setminus C_1$, its bit
would be one.

In the first case $\dist(q,C_0)\le e$. Project $q$ onto $C_0$, and
move the projection a further distance less than $e/2$ toward a fixed
interior point of $C_0$. This produces an interior solid start within
$3e/2$ of $q$, whose bit is zero. Uniform inballs permit a fixed small
convex-combination factor. The construction is measurable.

In the second case, when $q$ is outside the closed sweep, let $q_1$
be its metric projection onto $S_1(C_1)$ and set
$v=(q-q_1)/|q-q_1|$. At a boundary point choose an outward unit normal
from its compact normal cone, measurably, for example by the least
angle in a fixed angular convention. In both cases
$q\cdot v\ge p_{S_1(C_1)}(v)$, and support matching gives
$p_{S_0(C_0)}(v)\le p_{S_1(C_1)}(v)+e$. Hence
\[
                       q'=q+\tfrac32 e v
\]
lies strictly outside $S_0(C_0)$. It cannot enter another swept body
of table 0, by \eqref{eq:swept-separation} and smallness of $e$.
It is therefore a free miss, with bit zero. The displacement is $3e/2$.
Projection onto a closed convex set is continuous, so this map is
measurable on its case domain. Reverse the indices when the differing
bit is on side 1. Local finiteness and the separation give measurable
component selection. For $a=0$ every bit is zero and no change is needed.
Boundary cases can be assigned by the same inward or outward moves;
in particular the final forced-zero positions are not on boundaries.

The proof is pointwise in the nominal start and displacement. Apply it
to $(q,a)$ for a forward command and to $(q+a,-a)$ for its reverse.
Fresh nominal draws and deterministic perturbation maps preserve
conditional independence. They do not supply another observable.
\end{proof}

\begin{theorem}[Minimax resolution under table-dependent bounded errors]
\label{thm:calibration-floor}
For each fixed $s=6+\beta$ there is a fixed physical subclass of
$\mathcal B_\eta$, with known primitive lattice, two known species
and the fixed laboratory frame, and constants $c,r_0>0$ such that
for every $0<r<r_0$ two tables in this class have centered-support
$C^2$ distance at least
\begin{equation}\label{eq:calibration-separation}
                            c r^{(s-2)/s}
\end{equation}
yet admit position-error implementations bounded by $r$ with identical
laws for the entire adaptive bit experiment. This holds for every
nominal command design with $|a|\le t_{\max}<d_0$, arbitrary input
precision and arbitrary numbers of attempts. A sequential stopping
rule based on the bits cannot distinguish them either.

To specify the uncertainty quantifiers, let $\mathfrak I_r(\mathcal O)$
be the measurable implementation maps permitted above, and let
$\mathcal A_{t_{\max}}$ be the admissible bit-based controllers. The
construction has the stronger, controller-uniform form
\begin{gather}\label{eq:calibration-quantifiers}
 \forall\,0<r<r_0\quad
 \exists\,(\mathcal O_0,F_0),(\mathcal O_1,F_1),\quad
             F_j\in\mathfrak I_r(\mathcal O_j),\notag\\
 \forall\,\mathcal A\in\mathcal A_{t_{\max}}:\qquad
 \mathsf{Law}_{\mathcal O_0,F_0}^{\mathcal A}(Z)
       =\mathsf{Law}_{\mathcal O_1,F_1}^{\mathcal A}(Z).
\end{gather}
Here $Z$ includes the returned bits and the controller's recorded
commands and stopping decisions, not the actual launch positions.
The maps $F_0,F_1$ need not agree and may use the nominal draw and
command. They are fixed measurable maps for the two worlds, not an
extra observation or a common known jitter distribution.

Consequently, uniform success probability greater than $1/2$ at
centered $C^2$ accuracy $\nu$ under all such position errors requires
\begin{equation}\label{eq:calibration-necessary}
                            r\le C\nu^{s/(s-2)}.
\end{equation}
Constants, not an equality threshold, are asserted. The result is a
worst-case bounded-calibration statement; it does not assert this
obstruction for one fixed common offset or for a known mean-zero law.
\end{theorem}
\begin{proof}
Use the two-species periodic construction in Lemma~\ref{lem:packing},
but compare the unit-disk support $p_0=1$ with a single nonzero bump
\[
 p_1(\theta)=1+a h^s\psi((\theta-\theta_0)/h).
\]
The amplitude $a$ is fixed small enough that all curvature, separation
and patch margins remain uniform. Choose
$h=(r/(8a\|\psi\|_\infty))^{1/s}$. Then the Hausdorff distance of
the two variable bodies is at most $e=r/8$. The fixed disk and all
periodic copies are matched without alteration. At the bump center
the second derivatives differ by $a|\psi''(0)|h^{s-2}$; the
first-harmonic estimate in Lemma~\ref{lem:centered-packing} shows that
centering removes only $O(h^s)$. This gives
\eqref{eq:calibration-separation}. The lattice and disclosed data are
identical for both tables; the bump does not change primitiveness.

Lemma~\ref{lem:common-response} produces two physical implementations
with position errors at most $2e<r$. Couple the parameter-independent
controller randomness and every nominal draw across the two worlds.
Their first bits are equal by \eqref{eq:common-response}. Inductively,
identical histories generate identical next commands and identical
bits. Thus the complete bit streams, nominal controller records and
all bit-measurable stopping decisions have identical laws. No finite
or infinite sequence of those observations distinguishes the worlds.
An output law common to two disjoint accuracy sets cannot place
probability greater than $1/2$ in each. If $2\nu$ is smaller than
\eqref{eq:calibration-separation}, uniform success greater than $1/2$
is impossible. Rearranging proves \eqref{eq:calibration-necessary}.
\end{proof}

\subsection{The two resource requirements}
Together with Theorems~\ref{thm:adaptive} and \ref{thm:expected}, this
gives two necessary resources on one fixed nondegenerate subclass:
\[
 \sup_{\mathcal O}\E_{\mathcal O}\mathsf T
       \ge c\nu^{-1/(s-2)},\qquad
 r\le C\nu^{s/(s-2)}.
\]
Conversely the reciprocal compass construction attains the geometric
accuracy with the stated logarithmic overhead in attempts when
$\ell+\tau+t\alpha\le c\nu^{s/(s-2)}$. The converse calibration
necessity already holds with duration and angle exact, so it concerns
position error alone, not separate optimal exponents for all three
actuators. The lower bound allows any directions of length at most
$t_{\max}$, a larger design class than the four-atom upper construction.
The finite digital realization in Section~\ref{sec:digital} does not
alter either bound.

Thus the shrinking physical tolerance is not merely an omitted term in
a sample count: under the stipulated worst-case coupling it cannot be
replaced by additional samples. This is a statement within the active
collision model. Periodicity, the bounded presentation and the known
patch margin are still priors for the uniform global conclusion.
Neither crystallinity, apparatus calibration nor passive spectral
rigidity is inferred from the returned bits.

\section{Reconstruction with a fixed, unknown stationary launch law}
\label{sec:stationary}
The calibration converse quantifies over implementation maps that may
change with the table, the command and the nominal draw. Stationary
additive noise is a different experiment. We now give a constructive
inverse for this experiment at a fixed, nonzero launch spread. The
probability density need not be known, symmetric or mean zero. Its
support, and a lower bound on its mass near the support boundary, are
calibrated apparatus data.

\subsection{A common launch law and its geometric support}
Fix a known strictly convex body $K\subset\R^2$, containing the origin
in its interior, whose support function is $C^{6,\beta}$ and whose
curvature radii have positive lower and finite upper bounds. Assume
\begin{equation}\label{eq:fixed-footprint}
                       t+\diam K<d_0.
\end{equation}
This strict inequality, the bounds on $K$, and its position in the
laboratory frame are fixed independently of the target accuracy.
Let $\gamma\ge0$ and $b_0>0$ be fixed, with the following density class
nonempty:
\begin{equation}\label{eq:jitter-class}
 \mathcal J=\left\{j\in L^1(K):\ j\ge0,\quad\int_Kj=1,\quad
 j(z)\ge b_0\dist(z,\partial K)^\gamma
                  \text{ for almost every }z\in\operatorname{int}K\right\}.
\end{equation}
For $\gamma=0$ the lower bound is $b_0$ throughout the interior. No
upper bound, derivative bound or evaluation oracle for $j$ is needed.
The law $j$ is fixed during the experiment and is independent of the
commands, their compass direction and the table. The guarantee below
is uniform over this unknown nuisance law.

At nominal center $x$, draw fresh independent $Z\sim j$ and $a\sim\rho$.
The forward preparation attempts $(x+Z,a)$; its reciprocal reverse
attempts $(x+Z+a,-a)$ with the same prescribed joint distribution and
fresh draws. Only the pooled collision bit is returned. In particular,
$Z$, the actual start, free-start status and direction-resolved means
are not observed. Durations, directions and nominal centers are exact
in this stochastic model; certified rounding of the latter is handled
below. There is no additional arbitrary implementation error.

Set
\begin{equation}\label{eq:blurred-occupation}
 v_j(x)=\int_K j(z)\one_{\mathcal O}(x+z)\dd z,\qquad
 P_C=C+(-K).
\end{equation}
The first expression is a reflected convolution when $j$ is not even.
The same reciprocal identity gives $g_j=(T-I)v_j$. Every component
$P_C$ has diameter at most $D_0+\diam K$, and distinct such components
are at distance at least $d_0-\diam K>t$. Indeed, for $c\in C,d\in D$
and $z,w\in K$, one has
$|(c-z)-(d-w)|\ge |c-d|-|z-w|$.

\begin{lemma}[Support recovery without density deconvolution]
\label{lem:noise-support}
For every $j\in\mathcal J$,
\begin{equation}\label{eq:noise-support}
 \{v_j>0\}=\bigcup_C\operatorname{int}P_C,\qquad
             p_C^{\rm lab}=p_{P_C}^{\rm lab}-p_{-K}^{\rm lab}.
\end{equation}
Consequently the exact reciprocal mean functions on a sufficiently
large fixed protected aperture determine all protected component bodies,
without knowing $j$. With exact zero tests the full period group is then
recognized as in Lemma~\ref{lem:recognition}, without a supplied positive
nonperiod margin.
\end{lemma}
\begin{proof}
The intersection $C\cap(x+K)$ has nonempty interior exactly when
$x\in\operatorname{int}(C+(-K))$. One way to see the interior assertion
is to apply the open linear map $(c,z)\mapsto c-z$ to
$\operatorname{int}C\times\operatorname{int}K$; its image is the interior
of the Minkowski difference. On that intersection $j$ is positive
almost everywhere after translation, so its integral is positive.
Outside the interior, the intersection has area zero. The separation
above makes the positive components distinct. Support functions add
under Minkowski addition, which proves the second formula.

The proof of Lemma~\ref{lem:stopped} uses only $0\le v\le1$,
separated positive components of bounded diameter, and $g=(T-I)v$.
It does not require a smooth or known convolution kernel. Applied to
$v_j$, its killed iterates recover $v_j$ uniformly on the protected
region from its exact forcing. This proves the local assertion. The
period statement applies to the recovered original bodies, not to an
assumed lattice of the noise law. These support and Minkowski facts are
classical convex geometry; see \cite{Schneider}.
\end{proof}

\subsection{A uniform lower bound near the blurred boundary}
A quantitative version of positivity is needed for a finite experiment.
For a convex body $A$, write $A_{-r}=\{x:x+r\overline B\subset A\}$.
Whenever $r$ is smaller than its uniform interior rolling radius,
$p_{A_{-r}}=p_A-r$. The identity follows by subtracting $r$ from the
support function: its curvature measure remains nonnegative, and the
resulting Minkowski summand is exactly the erosion.

\begin{lemma}[Overlap and boundary mass]
\label{lem:cap-mass}
There are $c,e_0>0$, depending only on the physical class, $K$, $b_0$
and $\gamma$, such that for $0<e<e_0$ and every protected component,
\begin{equation}\label{eq:cap-mass}
 x\in(P_C)_{-e}\quad\Longrightarrow\quad
              v_j(x)\ge c e^{\gamma+3/2}.
\end{equation}
On the complement of $\bigcup_C\operatorname{int}P_C$ the value is zero.
Both assertions are uniform over $j\in\mathcal J$.
\end{lemma}
\begin{proof}
First consider the unweighted overlap
$w(x)=\area(C\cap(x+K))$. Write $P=C+(-K)$.
At $x_0\in\partial P$ with outer normal $n$, the uniquely determined
support points satisfy $x_0=y-z$, where $y\in\partial C$ has normal
$n$ and $z\in\partial K$ has normal $-n$. Let $r_C,r_K$ be fixed
positive interior rolling radii, smaller if necessary than their
uniform bounds. At $x=x_0-dn$, the intersection contains the lens of
the two disks with centers $y-r_Cn$ and $y+(r_K-d)n$ and radii
$r_C,r_K$. Their centers are separated by $r_C+r_K-d$.

Here is an explicit lower bound for the lens. In tangent-normal
coordinates about $y$, choose a fixed sufficiently small $c_1>0$.
For $|u|\le c_1\sqrt d$, the top of the first disk is
$-r_C+\sqrt{r_C^2-u^2}\ge-d/4$, whereas the bottom of the second is
$r_K-d-\sqrt{r_K^2-u^2}\le-3d/4$. For small $d$ the rectangle
\[
                  |u|\le c_1\sqrt d,\qquad
                         -3d/4\le y_n\le-d/4
\]
lies in both disks. Thus $w(x)\ge c_2d^{3/2}$ throughout a fixed
inner collar of $P$. An interior point at sufficiently small distance
$d$ from $\partial P$ has the displayed form at a nearest boundary
point, so this covers the whole collar, not only prescribed normals.

The bound extends to deeper points. The function $w^{1/2}$ is concave
on $P$: the Minkowski combination of the intersections at $x_1,x_2$
is contained in the intersection at their corresponding combination,
and the planar Brunn--Minkowski inequality gives concavity. Fix a small
collar depth $d_*$. On $\partial P_{-d_*}$ the preceding bound is
$c_2d_*^{3/2}$. A chord through any point of $P_{-d_*}$ has endpoints
on this boundary. Concavity bounds $w$ there by the same lower bound.
It follows that
\begin{equation}\label{eq:unweighted-cap}
 w(x)\ge c_3\min\{\dist(x,\partial P),d_*\}^{3/2},\qquad x\in P.
\end{equation}
The constants are uniform for the small erosions of $K$ used next:
their support norms and rolling radii remain uniformly controlled.

For $x\in P_{-e}$ put $K'=K_{-e/4}$. By the support identities,
$C+(-K')=P_{-e/4}$. Hence $x$ has distance at least $3e/4$ from
the boundary of this smaller sum. Applying \eqref{eq:unweighted-cap}
to $K'$ gives
$\area(C\cap(x+K'))\ge c_4e^{3/2}$. On $K'$ the density lower
bound is $b_0(e/4)^\gamma$. Restricting the defining integral to
this intersection proves \eqref{eq:cap-mass}. For $\gamma=0$ the
same argument works, although the erosion is unnecessary.
\end{proof}


\begin{proposition}[A density-independent inverse modulus]
\label{prop:jitter-modulus}
Consider two physical experiments $(\mathcal O_0,j_0)$ and
$(\mathcal O_1,j_1)$ with $j_0,j_1\in\mathcal J$ and the same calibrated
footprint. Let $g_0,g_1$ be their reciprocal pooled-mean differences and
put $\Delta=\|g_0-g_1\|_{W,\infty}$. On a common protected patch,
for sufficiently small $\Delta$, their complete components admit a
bijection satisfying
\begin{align}
 \max_C d_H(C_0,C_1)&\le C\Delta^{1/(\gamma+3/2)},
                                      \label{eq:jitter-Hausdorff}\\
 \max_C\|p_{C_0}^{\rm lab}-p_{C_1}^{\rm lab}\|_{C^2}
    &\le C\Delta^{(s-2)/(s(\gamma+3/2))}.
                                      \label{eq:jitter-modulus}
\end{align}
The densities need not agree. On $\mathcal B_\eta$, the primitive
orbit count and rational period relations in matched generating lists
also agree below the known
locking thresholds, and matched real period bases and free areas
have error bounded by the right side of \eqref{eq:jitter-Hausdorff}.
\end{proposition}
\begin{proof}
The comparison by stopping at the zero set of the first occupation in
Section~\ref{sec:queries}, and then interchanging the two occupations,
gives
$\|v_{j_0}-v_{j_1}\|_{U,\infty}\le H_K\Delta$.
This uses the bounded separated positive components $P_C$ and not
equality of the two densities. Put
$e=(2H_K\Delta/c)^{1/(\gamma+3/2)}$, with $c$ from
Lemma~\ref{lem:cap-mass}. For a point of $P_{C_0}$ choose a point of
$(P_{C_0})_{-e}$ at distance at most $e$; the erosion identity supplies
one. Its first occupation is at least $2H_K\Delta$, hence its second
occupation is positive. The two unions of expanded components are
therefore at Hausdorff distance at most $e$, by symmetry. Protective
collars include the test points for all components under comparison.

For $e$ below a fixed fraction of the expanded-component separation,
each connected expanded component lies in the $e$-neighborhood of a
unique component of the other union. The neighborhoods are pairwise
disjoint. Applying the same argument in reverse gives a bijection of
complete components. In support coordinates their distances are at
most $e$. Subtract the same function $p_{-K}$ to obtain
\eqref{eq:jitter-Hausdorff} for the original bodies.

For completeness, on a uniformly bounded $C^{6,\beta}$ class the
interpolation estimate
\[
 \|f\|_{C^2}\le C\|f\|_\infty^{(s-2)/s},\qquad
                              \|f\|_\infty\le1,
\]
follows by smoothing with a compact kernel of integral one and vanishing
moments through degree six. At smoothing scale $b$, its second
derivative is bounded by $C b^{-2}\|f\|_\infty$, and the derivative
remainder is $C b^{s-2}$. Choosing $b=\|f\|_\infty^{1/s}$ proves
the estimate; the first and zeroth derivatives have weaker bounds.
Apply this to the differences of matched support functions to prove
\eqref{eq:jitter-modulus}. Lemma~\ref{lem:locking} gives the final
assertion, with a matched independent-pair basis rather than a
continuously selected reduced lattice basis.
\end{proof}

\subsection{High-accuracy local queries at fixed launch spread}
The required query accuracy now tends to zero; a fixed $1/8$ error
would not locate the boundary of the positive set. A refinement of the
killed-aperture estimate avoids charging a spurious factor of the
iteration depth to every probability measurement.

\begin{lemma}[Aperture Green bound for perturbed iterates]
\label{lem:aperture-green}
Choose the fixed observation rectangle $W$ with collars protecting the
complete $P_C$ containing all targets in $U$. There are constants
$H_K,H_W,L_K$ such that the exact killed iterates for $g_j$ have error
at most $2^{-\lfloor n/L_K\rfloor}$ on $U$. If the forcing at each
state has error at most $\xi$ and each update error is at most $\zeta$,
the clipped approximate iterate satisfies
\begin{equation}\label{eq:jitter-green-error}
 \|\widehat V_n-v_j\|_{U,\infty}
       \le2^{-\lfloor n/L_K\rfloor}+H_W(\xi+\zeta).
\end{equation}
At any requested target $x$, all these iterates use only the states
$(x+t\Z^2)\cap W$, whose number is bounded by a fixed prior constant,
independently of $n$ and the target accuracy.
\end{lemma}
\begin{proof}
Take $H_K=(D_0+\diam K+t)^2/t^2$ and
$L_K=\max(1,\lceil2H_K\rceil)$. The stopped positive-component
argument proves the first assertion. For the walk killed on leaving
$W$, put $\tau_W=\inf\{n\ge0:X_n\notin W\}$ and take
$H_W=(\diam W+t)^2/t^2$. The bounded-stopping martingale calculation
gives $\sup_{x\in W}\E_x\tau_W\le H_W$.

Let $E_n$ be the pointwise absolute difference of the noisy and exact
killed iterates on the lattice through $x$. Averaging, positive part
and clipping to $[0,1]$ give
$E_{n+1}\le T^WE_n+\xi+\zeta$, with both errors zero outside $W$.
Iterating from zero yields
\[
 E_n(x)\le(\xi+\zeta)
       \sum_{k=0}^{n-1}(T^W)^k\one_W(x)
   =(\xi+\zeta)\E_x\min(n,\tau_W)\le H_W(\xi+\zeta).
\]
This bound uses exit from the entire fixed rectangle; it does not assume
that the approximate iterate knows the positive set. Combining it with
the exact error proves \eqref{eq:jitter-green-error}. A bounded rectangle
meets at most a fixed number of points of any translate of $t\Z^2$.
Transitions inside that state set are exact additions by the compass
step; transitions leaving it are killed, not wrapped or renormalized.
\end{proof}

\begin{proposition}[An indeterminate-layer query under fixed jitter]
\label{prop:jitter-query}
For each target $x\in U$, accuracy scale $0<e<e_0$ and conditional
failure allowance $0<\varepsilon<1/4$, a finite reciprocal experiment
returns the correct membership label of $\bigcup_CP_C$ whenever the
target lies outside its inner boundary layer of width $e$. In that
layer the label is unrestricted. It uses at most a fixed number of
command centers and
\begin{equation}\label{eq:jitter-query-cost}
                    C e^{-(2\gamma+3)}\log(C/\varepsilon)
\end{equation}
attempted collision bits. Its launch law has the same fixed support
$K$ at every query and every accuracy. No sample of the launch position
or of an occupation indicator is returned.
\end{proposition}
\begin{proof}
Let $b_e=c e^{\gamma+3/2}$ be the lower bound in
Lemma~\ref{lem:cap-mass}, decreasing $c$ so $b_e\le1$.
Estimate $v_j(x)$ to error at most $b_e/8$. In
Lemma~\ref{lem:aperture-green}, choose
$n=O(\log(C/b_e))$ with the first error at most $b_e/32$, and
estimate each pooled mean to error $b_e/(128H_W)$. Forcing is the
difference of two such means. Reserve update error $b_e/(64H_W)$.
These allocations give error less than $b_e/8$. There are only a
prior-bounded number of means, and their sample averages have Hoeffding
tails. A union bound gives the stated cost. Fresh preparations justify
the same assertion conditional on all previous data.

Threshold at $b_e/2$. Outside the positive components the true value is
zero, and at depth at least $e$ it is at least $b_e$. The test is therefore
correct at both kinds of points. Only the layer between them is
unrestricted. Rounding a target once by $O(e)$, followed by exact
addition of the dyadic compass step, enlarges this layer by $O(e)$;
it does not change the probability accuracy requirement. The finite
state iterates may be evaluated by rational arithmetic on the hit
counts or with the stated update reserve.
\end{proof}

\subsection{Uniform finite reconstruction and the role of calibration}
\begin{theorem}[Unknown stationary jitter on a fixed calibrated support]
\label{thm:stationary-jitter}
Fix the bounds of $\mathcal B_\eta$, a footprint $K$ satisfying
\eqref{eq:fixed-footprint}, and the nonempty density class
\eqref{eq:jitter-class}. There are constants $C,c,\nu_0>0$ and a finite
adaptive controller, independent of the unknown density $j$, such that
for $0<\nu<\nu_0$ and $0<\delta<1/4$,
\begin{equation}\label{eq:jitter-uniform-quantifier}
 \inf_{\mathcal O\in\mathcal B_\eta}\ \inf_{j\in\mathcal J}
 \Prob_{\mathcal O,j}\{\text{geometric loss}\le C\nu
              \text{ and the primitive discrete data are correct}\}
                         \ge1-\delta.
\end{equation}
The discrete data and the laboratory loss have the meanings fixed in
Section~\ref{sec:setting}. The controller uses one fixed launch support
$K$, the reciprocal compass law, and only the pooled collision bits.
It satisfies
\begin{align}
 J_\nu&\le C\nu^{-1/(s-2)}\log(C/\nu),\label{eq:jitter-centers}\\
 N_\nu&\le C\nu^{-((2\gamma+3)s+1)/(s-2)}
                  \log(C/\nu)\log(C/(\nu\delta)).\label{eq:jitter-attempts}
\end{align}
Every attempt, including a solid start or a miss, is counted. The
aperture is fixed by the priors. No localization scale of the random
start shrinks with $\nu$. For densities bounded below on $K$
($\gamma=0$), the power in \eqref{eq:jitter-attempts} is
$(3s+1)/(s-2)$.

Nominal centers have dyadic precision $O(\nu^{s/(s-2)})$ and require
$O(\log(C/\nu)+\log(C/\delta))$ digits per batch, including the
repetition count. Certified enclosures of the fixed numerical priors, including $b_0$
and $\gamma$, and certified support and derivative evaluations for $K$
are assumed when a numerical output is requested.
These are description and algorithmic assertions, not bounds on
manufacturing the launch law or calibrating its support.
\end{theorem}
\begin{proof}
The expanded bodies $P_C$ are strictly convex with uniformly bounded
$C^{6,\beta}$ support functions, positive curvature bounds, bounded
diameters and positive separation. Their curvature radii are the sums
of those of $C$ and $-K$. Choose a fixed enlarged protected aperture
for them and all coarse search brackets.

A fixed-scale version of Proposition~\ref{prop:jitter-query} gives
coarse hulls and interior centers for the expanded bodies. The proof of
Lemma~\ref{lem:coarse} applies verbatim to an indeterminate-layer label:
its deep grid nodes are retained and no node outside the expanded body
is retained; allowing an additional layer only increases the fixed
error reserve. The complete-component cutoff excludes fragments and
protects every required component. The isolated radial brackets and
positive inball therefore follow with changed prior constants.

Use $m$ angular directions, $h=2\pi/m\asymp\nu^{1/(s-2)}$, and
query at layer width $e=c h^s$. The relaxed bisection invariant of
Lemma~\ref{lem:bisection} needs only correctness outside a boundary
layer, not an exactly unbiased label or a known density. It returns
radial values of $P_C$ with error $O(e)$ in $O(\log(C/e))$ queries
per direction. Lemma~\ref{lem:radial} then gives smooth convex bodies
$\widehat P_C$ whose laboratory supports have $C^2$ error $O(h^{s-2})$
and supremum error $O(h^s)$.

Subtract the known footprint in support coordinates:
\begin{equation}\label{eq:estimated-subtraction}
                 \widehat p_C= p_{\widehat P_C}^{\rm lab}-p_{-K}^{\rm lab}.
\end{equation}
The true function $p_C+p_C''$ is bounded below by $1/\kappa_+$.
For small $\nu$, the $C^2$ error ensures
$\widehat p_C+\widehat p_C''>0$. Thus \eqref{eq:estimated-subtraction}
is the support function of an actual strictly convex body, not a formal
subtraction of sets. Its $C^2$ error is $O(h^{s-2})$ and its Hausdorff
error is $O(h^s)$. Subtracting a first harmonic for Steiner centering
preserves these orders. If a smooth output rather than a $C^{6,\beta}$
output is desired, use the finite smoothing described below.

Apply Lemma~\ref{lem:locking} to these original component bodies.
Their errors eventually lie below the known patch and bounded-denominator
thresholds, so the primitive orbit count and exact rational relations
are recovered. Real basis vectors and free area remain estimates, with
error $O(h^s)$. Positive separation and curvature yield a physical
periodic output exactly as in the proof of Theorem~\ref{thm:adaptive}.
The test is on the original bodies after subtraction; it is not an
assumption that an arbitrary noisy difference of sets preserves periods.

There are at most $Q_h=C(1+h^{-1}\log(C/h))$ adaptive local queries,
including coarse acquisition. Allocate conditional failure probability
$\delta/Q_h$ to each. Each query has a prior-bounded number of command
centers, even though its killed iteration depth grows logarithmically.
Equation~\eqref{eq:jitter-query-cost} yields total cost
\[
 C h^{-1}e^{-(2\gamma+3)}\log(C/h)\log(C/(h\delta)).
\]
Substituting $e=c h^s$ and $h\asymp\nu^{1/(s-2)}$ proves both counts.
Conditional confidence and the union bound prove
\eqref{eq:jitter-uniform-quantifier}. No property of $j$ other than its
support, lower bound and stationarity was used.

For numerical realization choose a dyadic compass step satisfying the
strict gap and one dyadic target mesh comparable to $e$. The state set
is the intersection of its translated exact compass lattice with a
fixed dyadic rectangle. Rounded radial bisection has the same mesh-floor
bound as in Theorem~\ref{thm:digital}. Sample counts here grow
polynomially in $1/\nu$, but their binary descriptions still require
only logarithmically many digits. The thresholds and count bounds can
be replaced by conservative dyadic enclosures of the fixed constants;
the receiver never needs to evaluate the unknown $j$.

For a finite smooth output, evaluate \eqref{eq:estimated-subtraction}
on an angular grid of width $b\asymp c\min(e^{1/3},\nu)$, to value
accuracy $O(b^3)$. The reconstructed supports have a uniform $C^3$
bound. Quadratic local interpolation patched with the fixed smooth
partition therefore has supremum error $O(b^3)$ and $C^2$ error
$O(b)$: Taylor remainders and value errors contribute the same orders
after up to two derivatives. These are respectively $O(e)$ and
$O(\nu)$. All support evaluations are computed from the finite radial
representation and the known footprint, not obtained through additional
collision observations. The curvature margin persists after reducing
constants. This final step is a disclosed numerical overhead, not part
of the attempted-bit count.
\end{proof}

\begin{corollary}[No fixed-jitter resolution floor in this experiment]
\label{cor:no-jitter-floor}
A fixed nondegenerate footprint satisfying \eqref{eq:fixed-footprint}
permits arbitrarily accurate reconstruction from finite collision bits,
uniformly over \eqref{eq:jitter-class}. In particular, stationary
uniform disk noise of any fixed radius $R$ with $t+2R<d_0$ has no
positive $C^2$ resolution floor in this model.
\end{corollary}
\begin{proof}
Keep $K$ and $j$ fixed and let $\nu$ decrease in
Theorem~\ref{thm:stationary-jitter}. Its sample bound is finite at every
positive accuracy. For a uniform disk, $\gamma=0$ and
$b_0=1/(\pi R^2)$ are admissible. The full density was not otherwise
needed by the construction.
\end{proof}

There is no contradiction with Theorem~\ref{thm:calibration-floor}.
The latter permits different, command- and draw-dependent
implementations in its indistinguishable worlds. Here all commands use
one stationary translate of a density with a calibrated support and
uniform boundary mass. The noise radius is not a worst-case actuator
error budget. An unknown displacement of the entire footprint, an
uncontrolled change of its support with the command, arbitrary errors
in duration or angle, and unbounded noise distributions are not included
in the new theorem. Nominal positioning and support calibration remain
separate apparatus requirements. Nor does the theorem claim a sharp
sample exponent: it proves that repetition can replace narrowing the
random launch spread in this stochastic experiment, not that every
kind of calibration can be replaced by repetition.

\section{Information categories and relation to earlier inverses}
\label{sec:comparison}
Active boundary estimation is an essential comparison for the finite
theorem, not merely an analogy with passive billiard spectra.
Castro and Nowak~\cite{CN} study adaptive binary classification under
boundary-fragment smoothness and noise assumptions, with matching
orders for their excess-risk criterion. Their input selects a feature
point whose response is a class label. This directly available label
is not our collision bit: a solid start and a free miss have the same
output here. Neither their observation model nor their loss identifies
our finite periodic table. We use an independent physical packing and
binary-tree proof rather than importing a classification lower bound.

Locatelli, Carpentier and Kpotufe~\cite{LCK} give an adaptive strategy
for smooth decision boundaries without prior knowledge of smoothness
and noise parameters. Their membership-query setting and local
line-search/interpolation construction make the closest algorithmic
comparison. In contrast, our smoothness is known, our error is a $C^2$
geometry norm, and the boundary label must first be computed from
reciprocal pooled collision means. We claim neither a new general
bisection principle nor an advance on adaptation to unknown smoothness.
The independent reduction in Proposition~\ref{prop:query} is what
permits this established active-estimation strategy in the physical
collision experiment.

The earlier scalar inverse reconstructed a uniformly accurate occupation
on a full two-dimensional grid. Its constructive upper bound was
$C\nu^{-3}\log(C/(\nu\delta))$ attempts on the same bounded
$C^{6,\beta}$ class, using $\sigma\asymp\nu^{3/2}$ and compensated
support smoothing. That theorem remains in the exact repository archive of A2 v31. The present
procedure uses $O(h^{-1})$ angular lines, a fixed-depth dependency
stencil for each queried point, and the full $C^{6,\beta}$ regularity.
It has a different adaptive command design, not a sharper analysis of
the old uniform design. We have not proved a matching lower bound
for that old design or for uniform spatial preparations.

The resource improvement is accompanied by the explicit localization
and calibration cost in \eqref{eq:precision}. Exact data remain a pair
of functionals on controlled spatial inputs. The lower bound concerns
only their finite binary realizations; it does not assert a low-information
passive invariant. The angle/time/position tolerances shrink with
localization and are not learned from the collision outcomes. Apparatus
motion, precision in representing input locations, and arithmetic
cost are distinct resources.

For completeness the exact v31 manuscript is preserved unchanged
inside the repository archive of the reviewed v33 source. It retains the arbitrary centered-law and general
nontrivial bounded-law inverses, the different-zero-set comparison,
the full fixed-aperture proof, rational interval enclosures, the
uniform-grid construction, and all earlier moment, marked-law and
count-fiber results in their original nested volumes. The more general
randomized law is an exact functional theorem; the finite local query
proved here uses the four-atom compass law. No continuous transition
law is discretized without an error budget. Rational intervals remain
conditional on forcing and survival bounds, not certificates of the
apparatus or physical prior.

The finite global conclusion also retains the known positive patch
margin and the bounded periodic presentation. Exact rational period
relations are distinct from approximate Euclidean generators. Removing
the known margin gives only pointwise eventual recognition, as described
in Section~\ref{sec:periods}. The result does not settle rigidity from
uniformly prepared count germs, exact analytic count germs, marked
length spectra or passive trajectories. These boundaries delimit the
information theorem without changing its topic or weakening any retained
mathematical assertion.

\subsection{Stationary noise versus adversarial implementation}
The new fixed-footprint theorem belongs also to support estimation with
contaminated data. Brunel, Klusowski and Yang~\cite{BKY} observe
continuous random vectors drawn from an unknown convex body and then
contaminated by additive Gaussian noise. Their geometric estimator
avoids Fourier inversion and is analyzed in Hausdorff loss. Here no
random vector or direct inside/outside response is observed: the data
are attempted collision bits at controlled centers. The fixed compact
noise support, its boundary-mass lower bound and the reciprocal reduction
are essential to our polynomial upper bound. The Gaussian observation
and its tail-based analysis are not covered by this theorem, and its
rates are not being compared as rates in a common experiment.

Support addition, inner rolling disks and the Brunn--Minkowski inequality
are classical \cite{Schneider}. Their role here is to remove an unknown
stationary density without learning that density. The positivity set of
the recovered blur, not its detailed values, identifies the footprint
sum. The high-accuracy query uses the Green bound for exit from a fixed
rectangle; its constants can be large and its finite-depth work grows
logarithmically. We do not claim these convex-geometric or random-walk
principles as new abstract theorems.

The two upper bounds cannot be ranked by their attempted-bit exponents
alone. The localized model attains the smaller sample power but requires
shrinking spread and worst-case tolerance. The stationary model pays a
larger, not claimed sharp sample power while keeping its random spread
fixed and allowing the density to be unknown. Its calibrated footprint,
boundary-mass condition, reciprocal stationarity and nominal positioning
are still apparatus assumptions. No theorem prices their manufacture,
extends the converse to every common noise law, or infers a passive
spectral invariant from these active experiments.

\subsection*{Source history and assisted analysis}
The article is self-contained and has no mandatory supplementary proof
volume. All proof chapters active in the reviewed v33 article remain
active here; the new stationary-noise results are additional conclusions.
The complete reviewed source and its earlier volumes are preserved
unchanged in the repository archive, with an explicit historical reading
map, rather than nested into the journal submission. They are not cited
as independently published articles. AI-assisted analysis contributed
to the new fixed-footprint inverse, boundary-mass estimates and experimental
comparisons. The named author is responsible for checking the proofs
and references before submission. Finite diagnostics and compilation
are not continuum proof certificates, apparatus validation or independent
referee acceptance.
\begin{thebibliography}{99}
\raggedright
\bibitem{CN}
R. M. Castro and R. D. Nowak, \emph{Minimax bounds for active learning},
IEEE Trans. Inform. Theory \textbf{54} (2008), 2339--2353.
\href{https://doi.org/10.1109/TIT.2008.920189}{doi:10.1109/TIT.2008.920189}.
\bibitem{LCK}
A. Locatelli, A. Carpentier, and S. Kpotufe,
\emph{An adaptive strategy for active learning with smooth decision boundary},
Proceedings of Algorithmic Learning Theory, Proceedings of Machine
Learning Research \textbf{83} (2018), 547--571.
\bibitem{LawlerLimic}
G. F. Lawler and V. Limic, \emph{Random Walk: A Modern Introduction},
Cambridge Studies in Advanced Mathematics \textbf{123}, Cambridge
University Press, 2010.
\bibitem{Shannon}
C. E. Shannon, \emph{A mathematical theory of communication},
Bell System Tech. J. \textbf{27} (1948), 379--423, 623--656.
\bibitem{Retained}
Q. Qi, \emph{Scalar collision laws and recognition of periodic dispersing
billiards}, A2 v31, October 4, 2026. Complete reviewed manuscript retained in the repository source archive;
not an independently published article or a required supplementary volume.
\bibitem{Schneider}
R. Schneider, \emph{Convex Bodies: The Brunn--Minkowski Theory},
second expanded edition, Encyclopedia of Mathematics and its Applications
\textbf{151}, Cambridge University Press, 2014.
\bibitem{BKY}
V.-E. Brunel, J. M. Klusowski, and D. Yang,
\emph{Estimation of convex supports from noisy measurements},
arXiv:1804.09879v1 (2018).
\end{thebibliography}

\end{document}
